% ============================================================================
% ProyectoFinal_LSA Salud.tex — Informe de Proyecto Final de Grado.
% Estructura: plantilla UTN-FRT 2023 + requisitos mínimos de la Ord. 1825
% (resumen de hasta ~200 palabras y palabras clave). Citas en APA.
% Compilar: latexmk -pdf -outdir=build "ProyectoFinal_LSA Salud.tex"
% ============================================================================
\documentclass[12pt, a4paper, oneside]{report}
\usepackage{lsasalud-pf}
\addbibresource{referencias.bib}

% Las figuras de Análisis y Diseño se reutilizan desde docs/TP7/img sin duplicarlas.
\graphicspath{{img/}{../TP7/img/}}
% \figuraTP{archivo sin extensión}{ancho}{leyenda}{etiqueta}
\newcommand{\figuraTP}[4]{%
    \begin{figure}[H]
        \centering
        \includegraphics[width=#2]{#1}
        \caption{#3}
        \label{#4}
    \end{figure}}

\hypersetup{pdftitle={\tituloProyecto}}

\begin{document}

\pagenumbering{roman}
\caratula
% Páginas preliminares: nota del tutor, agradecimientos, resumen y abstract.

\chapter*{Nota del tutor}
\addcontentsline{toc}{chapter}{Nota del tutor}
\pendiente{nota firmada por el tutor, dirigida al Director del Departamento de Sistemas (modelo en la plantilla UTN-FRT 2023), con apellido, nombre y legajo de cada autor. Se escanea y se inserta acá.}

\chapter*{Agradecimientos}
\addcontentsline{toc}{chapter}{Agradecimientos}
\pendiente{opcional (Ord. 1825, nota 1). Un párrafo por integrante.}

\chapter*{Resumen}
\addcontentsline{toc}{chapter}{Resumen}
\pendiente{hasta 200 palabras (Ord. 1825, 3.2): objetivos, metodología, resultados y conclusiones. Se escribe al final.}

\noindent\textbf{Palabras clave:} \pendiente{de 4 a 6 palabras clave (Ord. 1825, 3.3).}

\chapter*{Abstract}
\addcontentsline{toc}{chapter}{Abstract}
\pendiente{traducción al inglés del resumen (opcional, recomendado).}

\noindent\textbf{Keywords:} \pendiente{keywords.}


\tableofcontents
\listoffigures
\listoftables
\cleardoublepage
\pagenumbering{arabic}

\chapter{Resumen Ejecutivo}
\pendiente{una carilla: proyecto, objetivos, solución y conclusiones. Se escribe una vez terminado el informe.}

\chapter{Tema}
\pendiente{descripción del Sanatorio 9 de Julio, la necesidad (comunicación con pacientes sordos en la guardia) y los límites de la solución. Fuente: TP0 y TP1 (Acta de Constitución).}

\section{Alcance}
\pendiente{alcance del PMV y del sistema completo. Fuente: TP2 (Declaración de Alcance).}

\section{Fuera de alcance}
\pendiente{HCE, hardware a medida, capacitación en LSA. Fuente: TP2.}

\section{Supuestos y restricciones}
\pendiente{fuente: TP2.}

\chapter{Objetivos}

\section{Objetivo general}
\pendiente{propósito central, sin indicadores. Fuente: TP1 (Acta).}

\section{Objetivos específicos}
\pendiente{objetivos específicos del Acta: medibles, con plazo, realistas y verificables. Cada uno se retoma en el capítulo de Resultados.}

\chapter{Estado del Arte del Tema}
\pendiente{fundamentación teórica (LSA, reconocimiento de señas con visión por computadora, LSTM sobre landmarks) con citas académicas.}

\section{Contexto actual en Latinoamérica, Argentina y Tucumán}
\pendiente{población sorda, Ley 27.710 de reconocimiento de la LSA, situación en las guardias.}

\section{Innovación tecnológica del proyecto}
\pendiente{qué aporta el proyecto frente al estado del arte.}

\section{Impacto social del proyecto}
\pendiente{beneficiarios, de qué forma se benefician y otros impactos.}

\section{Necesidades y oportunidades que justifican el proyecto}
\pendiente{situación actual contra situación esperada.}

\section{Beneficios cuantitativos y cualitativos}
\pendiente{beneficios medibles y no medibles.}

\section{Tecnologías similares}
\pendiente{comparativa de traductores de lengua de señas existentes y justificación de la elección.}

\chapter{Gestión del proyecto}\label{cap:gestion}

Este capítulo resume cómo se planificó y gestionó el proyecto: las fases y los hitos, los criterios de éxito, el enfoque de trabajo, los interesados, los riesgos y el plan de gestión (estructura de trabajo, roles, comunicaciones y presupuesto). Las fuentes son los documentos de gestión elaborados durante el proyecto: Acta de Constitución, Registro de Interesados, Declaración de Alcance, Estructura de Desglose del Trabajo (EDT), Plan Financiero, Registro de Riesgos, Plan de Dirección e Informe Final.

\section{Fases del proyecto}

La Declaración de Alcance organiza el proyecto en seis fases. La tabla~\ref{tab:fases} las relaciona con los entregables y las fechas propuestas en dicho documento, ordenados cronológicamente.

\begin{table}[htbp]
\centering
\caption{Fases del proyecto, entregables y fechas propuestas.}\label{tab:fases}
\footnotesize
\begin{tabularx}{\textwidth}{@{}p{2.4cm}XXp{2.1cm}@{}}
\toprule
\textbf{Fase} & \textbf{Actividades principales} & \textbf{Entregables} & \textbf{Fecha propuesta} \\
\midrule
1. Relevamiento & Entrevistas presenciales con el personal administrativo y médico para identificar los requisitos del circuito de guardia. & E01: Documento de Alcance Validado & 30/05/2026 \\
2. Diseño & Arquitectura técnica (modelos de IA, lógica de \emph{backend} y \emph{frontend}) y selección de hardware comercial. & E02: Documento de Arquitectura & 15/06/2026 \\
3. Construcción del dataset & Recolección, filmación, clasificación y estructuración del corpus de señas clínicas, con validación lingüística de intérpretes de LSA. & E03: Dataset Clínico LSA & 30/06/2026 \\
4. Desarrollo y entrenamiento & Programación del PMV y entrenamiento de los modelos en condiciones controladas. & E04: Producto Mínimo Viable & 30/07/2026 \\
5. Pruebas e implementación piloto & Pruebas controladas y simuladas, seguidas del despliegue en un área crítica real del sanatorio. & E05: Informe de Pruebas Controladas; E06: Modelo de IA Optimizado; E07: Estación Desplegada en Guardia & 15/08/2026\newline 01/09/2026\newline 15/09/2026 \\
6. Evaluación y cierre & Medición de indicadores, optimización final, documentación técnica, manuales de uso y defensa académica. & E08: Informe de Impacto Operativo; E09: Sistema Final y Manuales; E10: Informe Final de Tesis & 15/10/2026\newline 01/11/2026\newline 15/11/2026 \\
\bottomrule
\end{tabularx}
\end{table}

El cronograma aprobado reprogramó algunas de estas fechas respecto de las propuestas originales: el diseño de la arquitectura se extendió hasta julio y las pruebas formales en entorno controlado se alinearon con el calendario de la cátedra (véase la sección de hitos). Para el seguimiento detallado, el trabajo se descompuso en paquetes de la EDT, presentados en el plan de gestión de este capítulo.

Las figuras~\ref{fig:gantt-fase1} a~\ref{fig:gantt-fase6} presentan el cronograma planificado (versión~4), una figura por fase de la EDT. Las barras oscuras agrupan cada entregable y las claras muestran sus paquetes de trabajo. Se trata de la línea base de planificación; las fechas reales de cada hito y su desvío se presentan en la sección siguiente.

% Generated from LSA Salud_Cronograma_V4.mpp (planned schedule). Do not edit by hand.
\begin{figure}[H]
\centering
\begin{ganttchart}[
  time slot format=isodate, x unit=6.000mm, y unit chart=0.62cm, y unit title=0.5cm,
  hgrid={draw=gray!25}, vgrid=false, title height=1, canvas/.append style={draw=gray!40},
  title label font=\scriptsize, bar label font=\scriptsize, group label font=\scriptsize\bfseries,
  bar label node/.append style={align=right, text width=6.4cm},
  group label node/.append style={align=right, text width=6.4cm},
  bar/.append style={fill=sectioncolor!45, draw=sectioncolor}, bar height=0.5,
  group/.append style={fill=ieeeblue, draw=ieeeblue}, group height=0.25,
  group left shift=0, group right shift=0, group peaks height=0.1
]{2026-05-12}{2026-05-20}
\gantttitlecalendar{month=shortname} \\
\ganttgroup{1.1 Entregable de documentación inicial}{2026-05-12}{2026-05-20} \\
\ganttbar{1.1.1 Acta de Constitución}{2026-05-12}{2026-05-15} \\
\ganttbar{1.1.2 Registro de interesados}{2026-05-15}{2026-05-20}
\end{ganttchart}
\caption{Cronograma planificado, fase 1 Documentación Inicial (12/05/2026 al 20/05/2026).}
\label{fig:gantt-fase1}
\end{figure}

% Generated from LSA Salud_Cronograma_V4.mpp (planned schedule). Do not edit by hand.
\begin{figure}[H]
\centering
\begin{ganttchart}[
  time slot format=isodate, x unit=2.024mm, y unit chart=0.62cm, y unit title=0.5cm,
  hgrid={draw=gray!25}, vgrid=false, title height=1, canvas/.append style={draw=gray!40},
  title label font=\scriptsize, bar label font=\scriptsize, group label font=\scriptsize\bfseries,
  bar label node/.append style={align=right, text width=6.4cm},
  group label node/.append style={align=right, text width=6.4cm},
  bar/.append style={fill=sectioncolor!45, draw=sectioncolor}, bar height=0.5,
  group/.append style={fill=ieeeblue, draw=ieeeblue}, group height=0.25,
  group left shift=0, group right shift=0, group peaks height=0.1
]{2026-05-20}{2026-06-30}
\gantttitlecalendar{month=shortname} \\
\ganttgroup{2.1 Entregable del alcance, matriz de trazabilidad y EDT}{2026-05-20}{2026-06-10} \\
\ganttbar{2.1.1 Matriz de trazabilidad de requisitos}{2026-05-20}{2026-05-25} \\
\ganttbar{2.1.2 Definición del alcance}{2026-05-22}{2026-05-26} \\
\ganttbar{2.1.3 EDT (Estructura de Desglose de Trabajo)}{2026-05-26}{2026-06-10} \\
\ganttgroup{2.2 Entregable de cronograma de trabajo del proyecto}{2026-06-10}{2026-06-19} \\
\ganttbar{2.2.1 Planificación de actividades}{2026-06-10}{2026-06-19} \\
\ganttbar{2.2.2 Confección del cronograma}{2026-06-12}{2026-06-18} \\
\ganttgroup{2.3 Entregable de viabilidad económica}{2026-06-16}{2026-06-30} \\
\ganttbar{2.3.1 Estimación económica del proyecto}{2026-06-16}{2026-06-23} \\
\ganttbar{2.3.2 Determinación de presupuesto}{2026-06-25}{2026-06-30} \\
\ganttgroup{2.4 Plan de gestión de riesgo}{2026-06-04}{2026-06-29} \\
\ganttbar{2.4.1 Identificar los riesgos}{2026-06-04}{2026-06-12} \\
\ganttbar{2.4.2 Planificar el plan de gestión de riesgos}{2026-06-22}{2026-06-29}
\end{ganttchart}
\caption{Cronograma planificado, fase 2 Plan de Gestión de Proyecto (20/05/2026 al 30/06/2026).}
\label{fig:gantt-fase2}
\end{figure}

% Generated from LSA Salud_Cronograma_V4.mpp (planned schedule). Do not edit by hand.
\begin{figure}[H]
\centering
\begin{ganttchart}[
  time slot format=isodate, x unit=0.934mm, y unit chart=0.62cm, y unit title=0.5cm,
  hgrid={draw=gray!25}, vgrid=false, title height=1, canvas/.append style={draw=gray!40},
  title label font=\scriptsize, bar label font=\scriptsize, group label font=\scriptsize\bfseries,
  bar label node/.append style={align=right, text width=6.4cm},
  group label node/.append style={align=right, text width=6.4cm},
  bar/.append style={fill=sectioncolor!45, draw=sectioncolor}, bar height=0.5,
  group/.append style={fill=ieeeblue, draw=ieeeblue}, group height=0.25,
  group left shift=0, group right shift=0, group peaks height=0.1
]{2026-05-26}{2026-08-24}
\gantttitlecalendar{month=shortname} \\
\ganttgroup{3.1 Informe de investigación}{2026-05-26}{2026-06-11} \\
\ganttbar{3.1.1 Obtención de tecnologías}{2026-05-26}{2026-06-04} \\
\ganttbar{3.1.2 Investigación de sistemas similares}{2026-06-04}{2026-06-11} \\
\ganttgroup{3.2 Documentación del análisis}{2026-06-11}{2026-08-24} \\
\ganttbar{3.2.1 Especificación de las funcionalidades de LSA Salud}{2026-06-11}{2026-06-22} \\
\ganttbar{3.2.2 Análisis del Dataset y modelo de IA}{2026-06-22}{2026-06-26} \\
\ganttbar{3.2.3 Análisis del módulo de texto a voz}{2026-06-22}{2026-06-25} \\
\ganttbar{3.2.4 Dataset médico especializado}{2026-07-03}{2026-07-21} \\
\ganttbar{3.2.5 Validación lingüística con intérpretes de LSA}{2026-07-21}{2026-08-24}
\end{ganttchart}
\caption{Cronograma planificado, fase 3 Análisis (26/05/2026 al 24/08/2026).}
\label{fig:gantt-fase3}
\end{figure}

% Generated from LSA Salud_Cronograma_V4.mpp (planned schedule). Do not edit by hand.
\begin{figure}[H]
\centering
\begin{ganttchart}[
  time slot format=isodate, x unit=2.931mm, y unit chart=0.62cm, y unit title=0.5cm,
  hgrid={draw=gray!25}, vgrid=false, title height=1, canvas/.append style={draw=gray!40},
  title label font=\scriptsize, bar label font=\scriptsize, group label font=\scriptsize\bfseries,
  bar label node/.append style={align=right, text width=6.4cm},
  group label node/.append style={align=right, text width=6.4cm},
  bar/.append style={fill=sectioncolor!45, draw=sectioncolor}, bar height=0.5,
  group/.append style={fill=ieeeblue, draw=ieeeblue}, group height=0.25,
  group left shift=0, group right shift=0, group peaks height=0.1
]{2026-06-26}{2026-07-24}
\gantttitlecalendar{month=shortname} \\
\ganttgroup{4.1 Diseño de arquitectura y APIs}{2026-06-26}{2026-07-15} \\
\ganttbar{4.1.1 Diseño de la arquitectura del sistema}{2026-06-26}{2026-07-03} \\
\ganttbar{4.1.2 Diseño de endpoints}{2026-07-03}{2026-07-09} \\
\ganttbar{4.1.3 Flujo de datos}{2026-07-09}{2026-07-15} \\
\ganttgroup{4.2 Diseño de interfaz de usuario}{2026-07-15}{2026-07-24} \\
\ganttbar{4.2.1 Wireframes de las pantallas}{2026-07-15}{2026-07-20} \\
\ganttbar{4.2.2 Maquetación estética y estilos (Mockups)}{2026-07-20}{2026-07-24}
\end{ganttchart}
\caption{Cronograma planificado, fase 4 Diseño (26/06/2026 al 24/07/2026).}
\label{fig:gantt-fase4}
\end{figure}

% Generated from LSA Salud_Cronograma_V4.mpp (planned schedule). Do not edit by hand.
\begin{figure}[H]
\centering
\begin{ganttchart}[
  time slot format=isodate, x unit=0.634mm, y unit chart=0.62cm, y unit title=0.5cm,
  hgrid={draw=gray!25}, vgrid=false, title height=1, canvas/.append style={draw=gray!40},
  title label font=\scriptsize, bar label font=\scriptsize, group label font=\scriptsize\bfseries,
  bar label node/.append style={align=right, text width=6.4cm},
  group label node/.append style={align=right, text width=6.4cm},
  bar/.append style={fill=sectioncolor!45, draw=sectioncolor}, bar height=0.5,
  group/.append style={fill=ieeeblue, draw=ieeeblue}, group height=0.25,
  group left shift=0, group right shift=0, group peaks height=0.1
]{2026-07-24}{2026-12-04}
\gantttitlecalendar{month=shortname} \\
\ganttgroup{5.1 Desarrollo del Frontend}{2026-07-24}{2026-09-21} \\
\ganttbar{5.1.1 Estructura y maquetación del cliente}{2026-07-24}{2026-09-14} \\
\ganttbar{5.1.2 Consumo de endpoints}{2026-09-14}{2026-09-21} \\
\ganttgroup{5.2 Desarrollo de la IA y procesamiento del Dataset}{2026-08-24}{2026-09-14} \\
\ganttbar{5.2.1 Preprocesamiento de datos}{2026-08-24}{2026-09-01} \\
\ganttbar{5.2.2 Entrenamiento y validación de la IA}{2026-09-01}{2026-09-14} \\
\ganttgroup{5.3 Desarrollo del backend}{2026-09-14}{2026-09-22} \\
\ganttbar{5.3.1 Servidor API de interconexión}{2026-09-14}{2026-09-22} \\
\ganttgroup{5.4 Pruebas del sistema}{2026-09-22}{2026-11-26} \\
\ganttbar{5.4.1 Pruebas E2E (End-To- End)}{2026-09-22}{2026-09-29} \\
\ganttbar{5.4.2 Pruebas con hardware físico}{2026-09-29}{2026-10-06} \\
\ganttbar{5.4.3 Validación con usuarios reales (Médicos)}{2026-10-06}{2026-11-26} \\
\ganttgroup{5.5 Presentación de versión final}{2026-11-26}{2026-12-04} \\
\ganttbar{5.5.1 Guía de usuario}{2026-11-26}{2026-12-01} \\
\ganttbar{5.5.2 Informe final del sistema}{2026-12-01}{2026-12-04}
\end{ganttchart}
\caption{Cronograma planificado, fase 5 Desarrollo del sistema (24/07/2026 al 04/12/2026).}
\label{fig:gantt-fase5}
\end{figure}

% Generated from LSA Salud_Cronograma_V4.mpp (planned schedule). Do not edit by hand.
\begin{figure}[H]
\centering
\begin{ganttchart}[
  time slot format=isodate, x unit=6.000mm, y unit chart=0.62cm, y unit title=0.5cm,
  hgrid={draw=gray!25}, vgrid=false, title height=1, canvas/.append style={draw=gray!40},
  title label font=\scriptsize, bar label font=\scriptsize, group label font=\scriptsize\bfseries,
  bar label node/.append style={align=right, text width=6.4cm},
  group label node/.append style={align=right, text width=6.4cm},
  bar/.append style={fill=sectioncolor!45, draw=sectioncolor}, bar height=0.5,
  group/.append style={fill=ieeeblue, draw=ieeeblue}, group height=0.25,
  group left shift=0, group right shift=0, group peaks height=0.1
]{2026-12-04}{2026-12-17}
\gantttitlecalendar{month=shortname} \\
\ganttgroup{6.1 Contenedorización y Empaquetado}{2026-12-04}{2026-12-14} \\
\ganttbar{6.1.1 Configuración de Entornos Aislados (Docker)}{2026-12-04}{2026-12-09} \\
\ganttbar{6.1.2 Creación de Scripts de Automatización}{2026-12-09}{2026-12-14} \\
\ganttgroup{6.2 Pruebas de despliegue}{2026-12-14}{2026-12-17} \\
\ganttbar{6.2.1 Pruebas de Funcionamiento en Entorno Desconectado}{2026-12-14}{2026-12-17}
\end{ganttchart}
\caption{Cronograma planificado, fase 6 Despliegue y seguridad de datos (04/12/2026 al 17/12/2026).}
\label{fig:gantt-fase6}
\end{figure}


\section{Hitos}

La tabla~\ref{tab:hitos} contrasta la fecha original de cada hito con la fecha real o pronosticada, según el seguimiento del Informe Final. Las filas con identificador \emph{H} son hitos del Acta de Constitución; el resto corresponde a entregables de la cátedra, cuya fecha original es el fin del período de seguimiento y cuya fecha real es la de la primera entrega. Una demora negativa indica entrega anticipada.

{\footnotesize
\begin{longtable}{@{}p{3.4cm}p{1.85cm}p{1.85cm}p{1.35cm}p{5.7cm}@{}}
\caption{Hitos y entregables: fechas originales, reales y demora (en días).}\label{tab:hitos}\\
\toprule
\textbf{Hito} & \textbf{Fecha original} & \textbf{Fecha real} & \textbf{Demora} & \textbf{Comentario} \\
\midrule
\endfirsthead
\toprule
\textbf{Hito} & \textbf{Fecha original} & \textbf{Fecha real} & \textbf{Demora} & \textbf{Comentario} \\
\midrule
\endhead
\midrule
\multicolumn{5}{r}{\emph{Continúa en la página siguiente}}\\
\endfoot
\bottomrule
\endlastfoot
Propuesta de tema (Traductor LSA Salud) & 29/04/2026 & 15/04/2026 & $-14$ & Propuestas subidas al inicio para recibir la devolución del cotutor; aprobada el 29/04. \\
H01. Presentación inicial del proyecto al cliente & 10/05/2026 & 10/05/2026 & 0 & Realizada en la fecha prevista, según confirmó el equipo. \\
Acta de Constitución e Interesados & 13/05/2026 & 03/05/2026 & $-10$ & Primera versión subida para pedir devolución y seguir iterando; aprobada el 13/05. \\
H02. Relevamiento de requerimientos & 20/05/2026 & 20/05/2026 & 0 & Realizado en la fecha prevista, según confirmó el equipo. \\
H03. Definición del alcance del proyecto & 30/05/2026 & 24/05/2026 & $-6$ & Alcance entregado en la primera versión de la Declaración de Alcance. \\
Declaración de Alcance & 03/06/2026 & 24/05/2026 & $-10$ & Primera versión el 24/05; EDT y matriz de trazabilidad en su versión 3 el 16/06; aprobada el 17/06. \\
H04. Diseño de arquitectura del sistema & 15/06/2026 & 06/07/2026 & $+21$ & Fecha del Acta, replanificada en el cronograma aprobado (diseño hasta julio). Tecnologías definidas el 06/07 según el repositorio; documento de arquitectura incluido en el entregable de Análisis y Diseño. \\
Cronograma & 24/06/2026 & 12/06/2026 & $-12$ & Versión 1 preliminar para recibir devolución (12/06); versión 4 el 25/06 y versión 5 el 29/10; aprobado el 25/06. \\
H05. Construcción del dataset clínico en LSA & 30/06/2026 & 02/07/2026 & $+2$ & Primer dataset propio con el modo de captura el 02/07 (historial del repositorio); se sigue ampliando. \\
Plan Financiero & 08/07/2026 & 07/07/2026 & $-1$ & Versión 1 el 07/07 y versión 3 el 11/07; estado resuelta, aprobado el 14/07. \\
H06. Desarrollo del PMV & 30/07/2026 & 17/07/2026 & $-13$ & Versión inicial funcional con reconocimiento básico y plantillas médicas, según el repositorio: el desarrollo empezó el 29/06 y quedó lista antes de la fecha. Primera entrega formal en el entregable de Implementación y Pruebas (21/10). \\
H07. Pruebas piloto en entorno controlado & 15/08/2026 & 20/10/2026 & $+66$ & Fecha del Acta, replanificada en el cronograma aprobado (25/06) y alineada al calendario de la cátedra. Ciclos formales de prueba el 20/10 y el 21/10. \\
Plan de Riesgos & 26/08/2026 & 02/09/2026 & $+7$ & Versión 1 el 02/09 y versión 2 el 08/09; aprobado el 14/09. El Plan de Dirección lo planifica para el 02/09. \\
H08. Ajustes y optimización del modelo & 01/09/2026 & 01/09/2026 & 0 & Realizado en la fecha prevista; optimización de latencia y ajuste del modelo en agosto. \\
H09. Implementación piloto en guardia & 15/09/2026 & --- & --- & Fuera del PMV de la materia; queda previsto para la etapa posterior del proyecto. \\
Plan de Gestión de Proyectos & 23/09/2026 & 18/09/2026 & $-5$ & Versión 1 el 18/09 y versión 2 el 25/09; aprobado el 04/10. \\
Análisis y Diseño (parcial) & 30/09/2026 & 28/09/2026 & $-2$ & Versión 1 el 28/09 y versión 2 el 18/10; aprobado el 28/10. \\
Análisis y Diseño (total) & 14/10/2026 & 28/09/2026 & $-16$ & Documentado en el entregable parcial por indicación del cotutor; cerrado el 28/10. \\
H10. Evaluación de impacto operativo & 15/10/2026 & --- & --- & Fuera del PMV; requiere el piloto en guardia. \\
H11. Presentación de versión final con manual de uso & 01/11/2026 & --- & --- & Fuera del PMV; se entregó el PMV ejecutable (28/10 y 02/11). \\
Implementación y Pruebas & 02/11/2026 & 21/10/2026 & $-12$ & Primera versión subida el 21/10, fecha del avance parcial; PMV el 28/10 y autoejecutable el 02/11; en evaluación del cotutor. \\
H12. Informe final de tesis y defensa & 15/11/2026 & 11/11/2026 & $-4$ & Exposición oral fijada por la cátedra para el 11/11/2026; se adelanta a la fecha del Acta. \\
Informe Final & 18/11/2026 & 18/11/2026 & 0 & Pronóstico igual a la fecha límite. \\
\midrule
\textbf{Total} & & & $\mathbf{-9}$ & \\
\end{longtable}
}

El balance acumulado es de nueve días de adelanto. Las demoras positivas son cuatro: el diseño de la arquitectura (21~días) y las pruebas piloto en entorno controlado (66~días), ambos por reprogramación del cronograma aprobado respecto de las fechas del Acta, además del Plan de Riesgos (7~días) y la construcción del dataset (2~días). El resto de las entregas se anticipó, en parte porque la fecha real registra la primera entrega y no la aprobación final. Los hitos de implementación piloto en guardia, evaluación de impacto operativo y presentación de la versión final con manual de uso quedan fuera del PMV y se prevén para la etapa posterior del proyecto.

\section{Resultados esperados e indicadores de evaluación}

Los criterios de aceptación del producto fijan los resultados esperados y la forma de medirlos. La tabla~\ref{tab:crit} los presenta junto con los objetivos y requisitos a los que responden.

\begin{table}[htbp]
\centering
\caption{Criterios de aceptación e indicadores de evaluación.}\label{tab:crit}
\footnotesize
\begin{tabularx}{\textwidth}{@{}lXXl@{}}
\toprule
\textbf{ID} & \textbf{Criterio de éxito} & \textbf{Indicador y medición} & \textbf{Relación} \\
\midrule
CRIT-01 & Latencia percibida inferior a 1~segundo bajo condiciones operativas normales. & Latencia promedio del sistema en condiciones normales. & OE07; RF-01; RNF-01 \\
CRIT-02 & \emph{Accuracy} mínimo del 90~\% en el reconocimiento de señas clínicas en entornos controlados. & \emph{Accuracy}, \emph{recall} y \emph{F1-score}. & OE06; RNF-02 \\
CRIT-03 & Precisión mínima del 85~\% en las pruebas operativas en la guardia real del sanatorio. & Precisión en pruebas operativas reales, con control de falsos positivos. & OE06; RF-05 \\
CRIT-04 & El 100~\% de las interacciones de video se procesa en tiempo real, sin almacenamiento persistente de imágenes o video en discos locales o externos. & Auditorías técnicas que verifiquen almacenamiento persistente igual a cero. & OE10; RNF-03 \\
CRIT-05 & El dataset clínico ampliado cuenta con el aval formal de la comunidad de intérpretes de LSA. & Aval formal de los intérpretes sobre el dataset. & OE04; RNF-05; RNF-06 \\
CRIT-06 & Las estaciones operan exclusivamente con componentes estándar de mercado (cámaras web comunes y pantallas duplicadas), sin hardware médico especializado. & Verificación de los componentes de las estaciones. & OE03; RNF-04; RNF-07 \\
CRIT-07 & Pantalla de consentimiento obligatoria antes del inicio de la consulta. & Presencia de la pantalla de consentimiento previa a la activación de la cámara. & OE10; RF-03 \\
\bottomrule
\end{tabularx}
\end{table}

Los criterios CRIT-01 a CRIT-06 figuran en la Declaración de Alcance; CRIT-07 se incorpora en la matriz de trazabilidad de requisitos y en el Plan de Dirección, asociado al requisito RF-03. Los diez objetivos específicos del capítulo anterior aportan las metas de impacto (tiempos, horas recuperadas, comprensión clínica, readmisiones, adopción y satisfacción), mientras que estos criterios definen las condiciones de aceptación técnica del producto. El contraste con los valores obtenidos se presenta en el capítulo de Resultados.

\section{Metodología de desarrollo}

El proyecto adoptó una metodología híbrida: \textbf{gestión del proyecto según la Guía del PMBOK del PMI, con desarrollo del producto iterativo e incremental}. La Guía del PMBOK aporta el marco de gestión de todo el ciclo, y el desarrollo iterativo e incremental organiza la construcción del software dentro de ese marco. Sus componentes son:
\begin{itemize}
  \item \textbf{Gestión según el PMBOK.} La iniciación, la planificación, el seguimiento y el cierre siguen los grupos de procesos de la Guía del PMBOK: acta de constitución, declaración de alcance, EDT con paquetes de trabajo que tienen responsable, esfuerzo y entradas y salidas, cronograma, plan financiero y registro de riesgos, con entregables asociados a fechas de control.
  \item \textbf{Construcción incremental del producto.} Primero se obtuvo una versión mínima funcional (el desarrollo comenzó el 29/06 y la versión inicial estuvo lista el 17/07), luego se ajustó el modelo y se optimizó la latencia (agosto) y finalmente se ejecutaron dos ciclos formales de prueba (20 y 21/10). El modo de captura permite ampliar el vocabulario de forma incremental.
  \item \textbf{Retroalimentación frecuente.} Varios entregables se presentaron en una primera versión para recibir la devolución del cotutor e iterar antes de su aprobación, y los documentos formales fueron revisados por el tutor y el cotutor.
\end{itemize}

La tabla~\ref{tab:metodologias} compara las alternativas consideradas para este contexto.

\begin{table}[htbp]
\centering
\caption{Comparación de enfoques de desarrollo frente a las condiciones del proyecto.}\label{tab:metodologias}
\footnotesize
\begin{tabularx}{\textwidth}{@{}p{2.6cm}XX@{}}
\toprule
\textbf{Enfoque} & \textbf{Fortalezas} & \textbf{Limitaciones en este proyecto} \\
\midrule
Cascada (secuencial) & Documentación y control claros; adecuado cuando los requisitos están cerrados. & Posterga la validación del modelo de IA hasta el final, cuando la precisión y la latencia son el principal riesgo técnico. \\
Iterativo e incremental & Entrega una versión funcional temprana y la mejora por ciclos; absorbe la incertidumbre técnica. & Requiere un marco de planificación y control para cumplir los entregables fijos de la cátedra. \\
Ágil (por ejemplo, Scrum) & Adaptación rápida y retroalimentación continua. & Presupone un equipo estable con ceremonias regulares; el equipo es de cuatro integrantes con carga académica variable (riesgos 22 y 25) y fecha final improrrogable. \\
\bottomrule
\end{tabularx}
\end{table}

\paragraph{Justificación.} Las condiciones del proyecto son una fecha límite improrrogable (15/11/2026), entregables formales con requisitos y criterios de aceptación trazados en una matriz, un equipo de cuatro integrantes y una incertidumbre técnica concentrada en el modelo de reconocimiento (riesgos 5, 9 y 16). Esas condiciones favorecen planificar y controlar el proyecto con la Guía del PMBOK y construir el producto de manera incremental, priorizando un PMV que permita medir latencia y precisión de forma temprana. El enfoque combinado conserva la trazabilidad exigida y reduce el riesgo técnico mejor que una secuencia estrictamente lineal.


\section{Registro de interesados}

El Registro de Interesados (versión del 05/06/2026) clasifica a cada parte según su postura (defensor, crítico o neutral) y la ubica en la matriz poder-interés para definir la estrategia de gestión. La tabla~\ref{tab:interesados} lo resume sin datos de contacto.

{\footnotesize
\begin{longtable}{@{}p{2.7cm}p{2.0cm}p{3.6cm}p{1.4cm}p{2.0cm}p{2.0cm}@{}}
\caption{Registro de interesados.}\label{tab:interesados}\\
\toprule
\textbf{Interesado} & \textbf{Rol} & \textbf{Expectativas principales} & \textbf{Tipo} & \textbf{Postura} & \textbf{Estrategia} \\
\midrule
\endfirsthead
\toprule
\textbf{Interesado} & \textbf{Rol} & \textbf{Expectativas principales} & \textbf{Tipo} & \textbf{Postura} & \textbf{Estrategia} \\
\midrule
\endhead
\midrule
\multicolumn{6}{r}{\emph{Continúa en la página siguiente}}\\
\endfoot
\bottomrule
\endlastfoot
Francisco Vicente & Tutor & Sistema robusto y alta precisión del traductor. & Externo & Defensor / crítico & Gestionar de cerca \\
Emmanuel López & Cotutor & No registradas. & Externo & Defensor / crítico & Gestionar de cerca \\
Sanatorio 9 de Julio & Cliente & Mejorar la comunicación y la eficiencia de la guardia. & Interno & Crítico & Gestionar de cerca \\
Estado (Ley 27.710) & Marco legal & No registradas. & Externo & Neutral & Mantener satisfecho \\
Recepcionista & Usuario & Comunicación directa y privada con el paciente. & Interno & Crítico & Gestionar de cerca \\
Enfermero/a & Usuario & Comunicación directa y privada con el paciente. & Interno & Crítico & Gestionar de cerca \\
Personas que se comunican con LSA & Usuario & Comunicación directa y privada con el personal. & Externo & Crítico & Gestionar de cerca \\
Natalia García Escudé & Doctora que aporta información & Que el grupo ponga en práctica el proyecto pronto. & Externo & Defensor & Mantener informado \\
Competencia & --- & Que el producto no supere al suyo. & Externo & Latente & --- \\
Creadores del dataset público de LSA & Aporte de datos & No registradas. & Externo & Neutral & Esfuerzo mínimo \\
Frameworks y librerías de código & MediaPipe, NumPy, Pandas, TensorFlow, OpenCV & No registradas. & Externo & Neutral & Esfuerzo mínimo \\
\end{longtable}
}

Los interesados de alto poder y alto interés (tutor, cotutor, sanatorio, recepcionista, enfermero/a y pacientes usuarios de LSA) se gestionan de cerca; el Estado y la competencia, de alto poder y bajo interés, se mantienen satisfechos o monitoreados; la doctora colaboradora, de bajo poder y alto interés, se mantiene informada; y los creadores del dataset y las librerías requieren esfuerzo mínimo. El Acta de Constitución y el Plan de Dirección identifican, además, a la comunidad de intérpretes de LSA (validación lingüística del dataset), a la cátedra y la universidad, a otras universidades referentes (UTN, UNC y UNLP) y al personal médico y administrativo como usuarios operativos del sistema.

\section{Gestión de riesgos del proyecto}

\subsection{Identificación y análisis cualitativo}

El Registro de Riesgos identifica 26 eventos, cada uno asociado a un paquete de la EDT. Para cada riesgo se estimó el impacto en una escala de 1 a 5, la probabilidad de ocurrencia y el área del proyecto afectada (tiempo, alcance o coste); el valor monetario esperado (VME) resulta del producto entre impacto y probabilidad, y la calificación clasifica el riesgo como inaceptable, aceptable con atención o aceptable. De los 26 riesgos, 13 son inaceptables, 11 son aceptables con atención y 2 son aceptables; todos figuran en estado abierto. La tabla~\ref{tab:riesgos} los detalla.

{\footnotesize
\setlength{\tabcolsep}{4pt}
\begin{longtable}{@{}p{0.6cm}p{0.95cm}p{5.4cm}p{2.0cm}p{0.9cm}p{1.0cm}p{0.9cm}p{2.1cm}@{}}
\caption{Registro de riesgos: identificación y análisis cualitativo.}\label{tab:riesgos}\\
\toprule
\textbf{ID} & \textbf{EDT} & \textbf{Descripción del riesgo} & \textbf{Área afectada} & \textbf{Imp.} & \textbf{Prob.} & \textbf{VME} & \textbf{Calificación} \\
\midrule
\endfirsthead
\toprule
\textbf{ID} & \textbf{EDT} & \textbf{Descripción del riesgo} & \textbf{Área afectada} & \textbf{Imp.} & \textbf{Prob.} & \textbf{VME} & \textbf{Calificación} \\
\midrule
\endhead
\midrule
\multicolumn{8}{r}{\emph{Continúa en la página siguiente}}\\
\endfoot
\bottomrule
\endlastfoot
1 & 3.2.5 & Falta de disponibilidad de intérpretes oficiales de LSA para validar el dataset clínico y las señas médicas. & Tiempo, alcance, coste & 5 & 30~\% & 1,5 & Inaceptable \\
2 & 5.2.2 & Sobreajuste del modelo durante el entrenamiento, que limita la generalización de las señas. & Tiempo, alcance & 4 & 40~\% & 1,6 & Aceptable con atención \\
3 & 1.1.2 & Resistencia o falta de colaboración del personal de guardia en la recolección de requisitos. & Tiempo & 3 & 50~\% & 1,5 & Aceptable con atención \\
4 & 6.2.1 & Iluminación deficiente en los boxes de admisión que impida a la cámara captar rasgos y profundidad. & Alcance & 5 & 25~\% & 1,25 & Inaceptable \\
5 & 4.1.3 & Latencia del procesamiento de video con IA superior al umbral de 1~segundo. & Alcance, coste & 5 & 35~\% & 1,75 & Inaceptable \\
6 & 5.2.1 & Ruido visual del entorno hospitalario que afecte la extracción de coordenadas (x, y, z) de la postura. & Alcance, tiempo & 4 & 60~\% & 2,4 & Aceptable con atención \\
7 & 6.2.1 & Interrupción del suministro eléctrico o de la red local durante las pruebas piloto. & Tiempo & 3 & 45~\% & 1,35 & Aceptable con atención \\
8 & 6.1.1 & Hardware local del sanatorio con capacidad insuficiente para ejecutar con fluidez los contenedores Docker. & Tiempo, coste & 4 & 20~\% & 0,8 & Inaceptable \\
9 & 3.2.4 & Dificultad para compilar y grabar suficientes muestras de LSA de vocabulario médico específico (por ejemplo, niveles de dolor). & Alcance, tiempo & 5 & 30~\% & 1,5 & Inaceptable \\
10 & 4.2.1 & Diseño de interfaz confuso o poco ergonómico para pacientes sordos durante la atención. & Tiempo, alcance & 3 & 40~\% & 1,2 & Aceptable con atención \\
11 & 3.2.3 & Audio poco natural o incomprensible del módulo de texto a voz sin conexión. & Alcance & 2 & 50~\% & 1,0 & Aceptable \\
12 & 5.4.2 & Falta de stock local o aumento de precios de las pantallas de doble cara y cámaras web estándar. & Tiempo, coste & 4 & 35~\% & 1,4 & Inaceptable \\
13 & 4.1.1 & Restricciones legales o cambios en normativas de privacidad del sanatorio sobre la captura de video de pacientes. & Alcance, tiempo, coste & 5 & 15~\% & 0,75 & Inaceptable \\
14 & 2.2.1 & Renuncia, enfermedad o reasignación de un desarrollador clave del equipo. & Tiempo, coste & 5 & 20~\% & 1,0 & Inaceptable \\
15 & 5.4.3 & El personal administrativo prefiere su método tradicional y no se adapta a la estación fija. & Alcance, tiempo & 4 & 45~\% & 1,8 & Aceptable con atención \\
16 & 5.2.2 & Falla del modelo de IA al capturar e interpretar gestos faciales no manuales vitales en la sintaxis de la LSA. & Alcance & 5 & 30~\% & 1,5 & Inaceptable \\
17 & 2.3.2 & Inflación e inestabilidad económica que superen la reserva de gestión y contingencia del presupuesto. & Coste & 3 & 70~\% & 2,1 & Aceptable con atención \\
18 & 5.4.1 & Problemas de integración y sincronización de datos entre el \emph{frontend} y el motor de inferencia. & Tiempo & 3 & 50~\% & 1,5 & Aceptable con atención \\
19 & 5.2.1 & Fallo del algoritmo al normalizar el tamaño de manos y cuerpo según la distancia del paciente a la pantalla. & Alcance, tiempo & 4 & 40~\% & 1,6 & Aceptable con atención \\
20 & 3.1.2 & Publicación de una tecnología abierta de traducción médica que vuelva obsoleta parte de la base tecnológica. & Alcance & 2 & 15~\% & 0,3 & Aceptable \\
21 & 5.4.3 & El sanatorio revoca o posterga el acceso físico a los boxes de triaje para el piloto. & Tiempo, coste & 5 & 20~\% & 1,0 & Inaceptable \\
22 & 5.1.1 & Un integrante reduce drásticamente su participación por no regularizar materias correlativas. & Alcance, tiempo & 5 & 30~\% & 1,5 & Inaceptable \\
23 & 5.3.1 & Retraso severo en tareas de la ruta crítica con holgura cero, que compromete la fecha de entrega de la tesis. & Tiempo, coste & 5 & 35~\% & 1,75 & Inaceptable \\
24 & 5.4.1 & Flujos de datos entre \emph{frontend} y modelo de IA con latencias inmanejables, que exigen reescribir la lógica de conexión. & Alcance, tiempo, coste & 5 & 30~\% & 1,5 & Inaceptable \\
25 & 5.1.1 & Un integrante disminuye su ritmo durante los exámenes parciales, sin abandonar el proyecto. & Tiempo & 3 & 70~\% & 2,1 & Aceptable con atención \\
26 & 5.2.2 & El entrenamiento toma un 15~\% más de cómputo del planificado para superar la precisión mínima del 85~\% (CRIT-03). & Tiempo, coste & 3 & 50~\% & 1,5 & Aceptable con atención \\
\end{longtable}
}

La tabla~\ref{tab:matriz} ubica los riesgos en la matriz de impacto y probabilidad; los identificadores en negrita corresponden a riesgos inaceptables. Los 11 riesgos de impacto 5 son todos inaceptables, entre ellos la disponibilidad de intérpretes, la latencia, la disponibilidad de muestras clínicas, la captura de gestos no manuales y la ruta crítica.

\begin{table}[htbp]
\centering
\caption{Matriz de impacto y probabilidad (identificadores de riesgo).}\label{tab:matriz}
\small
\begin{tabularx}{\textwidth}{@{}lXXXX@{}}
\toprule
\textbf{Probabilidad} & \textbf{Impacto 2} & \textbf{Impacto 3} & \textbf{Impacto 4} & \textbf{Impacto 5} \\
\midrule
15~\% & 20 & --- & --- & \textbf{13} \\
20~\% & --- & --- & \textbf{8} & \textbf{14}, \textbf{21} \\
25~\% & --- & --- & --- & \textbf{4} \\
30~\% & --- & --- & --- & \textbf{1}, \textbf{9}, \textbf{16}, \textbf{22}, \textbf{24} \\
35~\% & --- & --- & \textbf{12} & \textbf{5}, \textbf{23} \\
40~\% & --- & 10 & 2, 19 & --- \\
45~\% & --- & 7 & 15 & --- \\
50~\% & 11 & 3, 18, 26 & --- & --- \\
60~\% & --- & --- & 6 & --- \\
70~\% & --- & 17, 25 & --- & --- \\
\bottomrule
\end{tabularx}
\end{table}

\subsection{Plan de respuesta}

Los 13 riesgos inaceptables cuentan con un plan de respuesta que combina acciones para evitar, mitigar o transferir el riesgo y acciones de contingencia, estas últimas con cargo, cuando corresponde, a la reserva de gestión. La tabla~\ref{tab:respuesta} lo sintetiza.

{\footnotesize
\setlength{\tabcolsep}{4pt}
\begin{longtable}{@{}p{1.5cm}p{3.5cm}p{5.0cm}p{4.9cm}@{}}
\caption{Plan de respuesta a los riesgos inaceptables.}\label{tab:respuesta}\\
\toprule
\textbf{Riesgo} & \textbf{Evitar o transferir} & \textbf{Mitigar} & \textbf{Contingencia} \\
\midrule
\endfirsthead
\toprule
\textbf{Riesgo} & \textbf{Evitar o transferir} & \textbf{Mitigar} & \textbf{Contingencia} \\
\midrule
\endhead
\midrule
\multicolumn{4}{r}{\emph{Continúa en la página siguiente}}\\
\endfoot
\bottomrule
\endlastfoot
1. Intérpretes de LSA & Validación asincrónica con secuencias grabadas de 30 \emph{frames}. & Contacto temprano con la Confederación Argentina de Sordos o asociaciones locales; dos integrantes depuran y estructuran el dataset antes de enviarlo. & Entrenar el PMV solo con LSA64 de manera temporal; contratar intérpretes externos por videollamada con la reserva de gestión. \\
4. Iluminación & Instalar las estaciones solo en áreas con iluminación comprobada. & Aumento de datos; pruebas sobre hardware físico; cámaras de 720p/1080p con buena compensación de exposición. & Paneles LED direccionales con la reserva de contingencia; reubicar el piloto en otro box. \\
5. Latencia & Extraer solo las coordenadas (x, y, z) con MediaPipe en lugar de procesar píxeles. & Inferencia 100~\% local; modelo espacio-temporal que clasifica solo vectores numéricos; medir los tiempos durante el desarrollo del servidor API. & Procesar solo los 21 puntos de manos y rostro; dotar a la estación de una GPU básica con la reserva. \\
8. Hardware del sanatorio & Restringir el diseño a componentes estándar de mercado. & Solo coordenadas (más liviano que el video); scripts de inicio y parada; pruebas tempranas sobre el hardware físico. & GPU básica con la reserva; trasladar la inferencia a un servidor local de la intranet. \\
9. Muestras de vocabulario médico & Reuniones previas con intérpretes para confirmar el vocabulario antes de grabar. & Prototipar con LSA64; grabar coordenadas en lugar de video. & Aumento de datos propio; grabación piloto de emergencia reducida a 10 señas críticas con 5 personas. \\
12. Stock y precios de hardware & Hardware comercial genérico; \emph{transferir:} comprar los componentes en las etapas iniciales. & Monitorear inflación y divisas; cotizar en dólares a proveedores e inversores. & Reserva de gestión para sobrecostos; soporte propio que una dos monitores o tabletas. \\
13. Normativa de privacidad & Arquitectura sin almacenamiento persistente (CRIT-04); inferencia 100~\% local. & Extraer solo características numéricas; consentimiento explícito antes de capturar señas. & Limitar las pruebas a laboratorio hasta la aprobación legal del sanatorio; operar solo con seguimiento de manos. \\
14. Desarrollador clave & Confirmar tempranamente los compromisos académicos. & Revisiones de código cruzadas; validación grupal de las estimaciones de la EDT; documentación técnica exhaustiva. & Contratar un recurso temporal con la reserva; reasignar tareas al equipo. \\
16. Gestos no manuales & Variar distancia e iluminación en la captura propia de datos. & Aumento de datos sintético; evaluación continua con datos de validación; limitar el vocabulario al triaje de baja urgencia. & Ajustar pesos según las matrices de confusión; reducir el PMV a 10 señas críticas de mayor confianza. \\
21. Acceso al sanatorio & Formalizar acuerdos tempranos de disponibilidad de las instalaciones. & Docker y scripts para portabilidad; hardware que no requiera intervenir la red; pruebas E2E en un laboratorio de la facultad. & Trasladar el piloto al Hospital Nuestra Señora del Carmen; replantear el cronograma con la reserva financiera. \\
22. Participación reducida & Consultar la carga académica real antes de asignar paquetes largos. & Reuniones grupales periódicas para redistribuir cargas; documentación técnica exhaustiva. & Reasignar tareas; apoyo temporal con la reserva. \\
23. Ruta crítica & Lecciones aprendidas y juicio de expertos al estimar. & Coeficientes de contingencia de tiempo; monitoreo estricto de las tareas sin holgura. & Compresión (\emph{crashing}) con las reservas; desarrollo paralelo de \emph{frontend} y \emph{backend}. \\
24. Flujos \emph{frontend}--modelo & Definir el contrato de mensajes antes de programar: el \emph{frontend} envía solo coordenadas por una conexión WebSocket persistente, sin imágenes ni una solicitud HTTP por cuadro. & Enviar ventanas de cuadros con un paso de predicción fijo; prototipo temprano de la conexión de punta a punta; medir la latencia contra CRIT-01 en cada iteración. & Aumentar el paso entre predicciones o acortar la ventana; ejecutar el \emph{frontend} y el servidor de inferencia en el mismo equipo de la estación. \\
\end{longtable}
}

Varias de las medidas de evitación y mitigación se materializaron en el PMV: el procesamiento es íntegramente local, el \emph{backend} recibe solo las coordenadas de las manos y la estación del paciente exige consentimiento antes de activar la cámara, lo que responde a los riesgos 5, 8 y 13.

\section{Plan de gestión del proyecto}

\subsection{Estructura de desglose del trabajo}

La EDT descompone el proyecto en seis fases, dieciséis entregables y 36 paquetes de trabajo, con un esfuerzo total de 1.136~horas. La tabla~\ref{tab:edt} la resume hasta el nivel de paquete.

{\footnotesize
\begin{longtable}{@{}p{3.0cm}p{3.4cm}p{6.9cm}p{1.1cm}@{}}
\caption{Estructura de desglose del trabajo: fases, entregables y paquetes.}\label{tab:edt}\\
\toprule
\textbf{Fase (nivel 1)} & \textbf{Entregable (nivel 2)} & \textbf{Paquetes de trabajo (nivel 3)} & \textbf{Horas} \\
\midrule
\endfirsthead
\toprule
\textbf{Fase (nivel 1)} & \textbf{Entregable (nivel 2)} & \textbf{Paquetes de trabajo (nivel 3)} & \textbf{Horas} \\
\midrule
\endhead
\midrule
\multicolumn{4}{r}{\emph{Continúa en la página siguiente}}\\
\endfoot
\bottomrule
\endlastfoot
\textbf{1. Documentación inicial} (24~h) & 1.1 Documentación inicial & 1.1.1 Acta de Constitución; 1.1.2 Registro de interesados & 24 \\
\midrule
\textbf{2. Plan de gestión del proyecto} (216~h) & 2.1 Alcance, matriz de trazabilidad y EDT & 2.1.1 Matriz de trazabilidad de requisitos; 2.1.2 Definición del alcance; 2.1.3 EDT & 68 \\
 & 2.2 Cronograma & 2.2.1 Planificación de actividades; 2.2.2 Confección del cronograma & 48 \\
 & 2.3 Viabilidad económica & 2.3.1 Estimación económica; 2.3.2 Determinación del presupuesto & 44 \\
 & 2.4 Plan de gestión de riesgos & 2.4.1 Identificar los riesgos; 2.4.2 Planificar la gestión de riesgos & 56 \\
\midrule
\textbf{3. Análisis} (250~h) & 3.1 Informe de investigación & 3.1.1 Obtención de tecnologías; 3.1.2 Investigación de sistemas similares & 68 \\
 & 3.2 Documentación del análisis & 3.2.1 Especificación de funcionalidades; 3.2.2 Análisis del dataset y del modelo de IA; 3.2.3 Análisis del módulo de texto a voz; 3.2.4 Dataset médico especializado; 3.2.5 Validación lingüística con intérpretes de LSA & 182 \\
\midrule
\textbf{4. Diseño} (110~h) & 4.1 Diseño de arquitectura y APIs & 4.1.1 Arquitectura del sistema; 4.1.2 Endpoints; 4.1.3 Flujo de datos & 72 \\
 & 4.2 Diseño de interfaz de usuario & 4.2.1 Wireframes; 4.2.2 Maquetación estética y estilos & 38 \\
\midrule
\textbf{5. Desarrollo del sistema} (484~h) & 5.1 Desarrollo del \emph{frontend} & 5.1.1 Estructura y maquetación del cliente; 5.1.2 Consumo de endpoints & 236 \\
 & 5.2 Desarrollo de la IA y procesamiento del dataset & 5.2.1 Preprocesamiento de datos; 5.2.2 Entrenamiento y validación de la IA & 88 \\
 & 5.3 Desarrollo del \emph{backend} & 5.3.1 Servidor API de interconexión & 36 \\
 & 5.4 Pruebas del sistema & 5.4.1 Pruebas E2E; 5.4.2 Pruebas con hardware físico; 5.4.3 Validación con usuarios reales (médicos) & 88 \\
 & 5.5 Presentación de la versión final & 5.5.1 Guía de usuario; 5.5.2 Informe final del sistema & 36 \\
\midrule
\textbf{6. Despliegue y seguridad de datos} (52~h) & 6.1 Contenedorización y empaquetado & 6.1.1 Entornos aislados (Docker); 6.1.2 Scripts de automatización & 30 \\
 & 6.2 Pruebas de despliegue & 6.2.1 Pruebas de funcionamiento en entorno desconectado & 22 \\
\end{longtable}
}

\subsection{Roles y responsabilidades}

La tabla~\ref{tab:roles} presenta la estructura del equipo. Cada paquete de la EDT tiene un integrante responsable y la matriz de asignación de responsabilidades del Plan de Dirección indica quién responde y quién participa en cada uno. Medido por el esfuerzo de los paquetes que cada integrante lidera, la carga es de 248 horas para Hamada, 300 para Arnedo, 296 para García Tarcic y 292 para Véliz, lo que suma las 1.136 horas de la EDT.

\begin{table}[htbp]
\centering
\caption{Roles y responsabilidades.}\label{tab:roles}
\footnotesize
\begin{tabularx}{\textwidth}{@{}p{3.1cm}p{3.3cm}X@{}}
\toprule
\textbf{Integrante} & \textbf{Rol} & \textbf{Responsabilidades} \\
\midrule
Hamada, Santino & Director de Proyecto y Líder Técnico & Coordinación general, gestión del cliente, planificación, control del cronograma, asignación de tareas, definición de la arquitectura y supervisión técnica. \\
García Tarcic, Iñaki Ezequiel & Analista Funcional y Desarrollador Backend & Relevamiento de requisitos, definición de KPI y evaluación de impacto; diseño e implementación de la lógica del servidor, las APIs y la integración de servicios. \\
Véliz, Héctor Matías & QA y Testing; Desarrollador Frontend & Diseño de pruebas, validación del sistema y medición de métricas operativas; diseño y desarrollo de la interfaz de usuario. \\
Arnedo, Emmanuel Augusto & Desarrollador IA; Fullstack y QA & Implementación de los modelos de visión por computadora y optimización del sistema; desarrollo integral de \emph{frontend} y \emph{backend} y aseguramiento de la calidad. \\
Vicente, Francisco José & Tutor & Supervisión académica y técnica principal; orientación metodológica y directrices estratégicas. \\
López, Emmanuel & Cotutor & Asistencia en la tutoría y supervisión técnica; acompañamiento en las decisiones de diseño y ejecución. \\
\bottomrule
\end{tabularx}
\end{table}

Como subcontratistas y proveedores, el Plan de Dirección identifica a los intérpretes de LSA (validación lingüística del dataset y asesoramiento en vocabulario clínico), los proveedores de hardware (cámaras y pantallas) y las instituciones de salud locales (espacio para pruebas piloto y validación en entorno real).

\subsection{Comunicaciones}

La documentación del proyecto no incluye un plan de comunicaciones formal; las prácticas registradas son las siguientes. El tutor, el cotutor, el sanatorio y los usuarios (recepción, enfermería y pacientes) se gestionan de cerca según el Registro de Interesados. El relevamiento se realizó mediante entrevistas con el personal de salud y una reunión con el cliente para definir los requisitos. Los documentos formales se redactaron en conjunto y fueron revisados por el tutor y el cotutor, y varios entregables se subieron en una primera versión para recibir su devolución. Como medida de gestión de riesgos, se previeron reuniones grupales periódicas para evaluar la carga del equipo y redistribuir tareas.

Los canales utilizados fueron los siguientes:
\begin{itemize}
  \item \textbf{Comunicación interna del equipo:} informal, por WhatsApp para la coordinación diaria y por Discord para reuniones de trabajo con voz y pantalla compartida. Las reuniones no tuvieron una frecuencia fija y se convocaron según las necesidades de cada entregable.
  \item \textbf{Comunicación con la cátedra:} formal, a través de Redmine, donde se registraron las tareas de cada entregable, se subieron los documentos y se recibieron las devoluciones del tutor y del cotutor.
  \item \textbf{Trabajo técnico:} el código y la documentación se versionaron en un repositorio Git, con cambios integrados mediante \emph{pull requests}.
\end{itemize}

\subsection{Presupuesto}

El Plan Financiero estima la ejecución entre el 12/05/2026 y el 17/12/2026, con un esfuerzo total de 1.136~horas. Los importes se expresan en pesos argentinos; la tabla~\ref{tab:presupuesto} los resume por rubro, con y sin la contingencia asignada a cada uno.

\begin{table}[htbp]
\centering
\caption{Resumen del presupuesto por rubro (en pesos, sin IVA salvo indicación).}\label{tab:presupuesto}
\footnotesize
\begin{tabularx}{\textwidth}{@{}p{3.0cm}Xrr@{}}
\toprule
\textbf{Rubro} & \textbf{Detalle} & \textbf{Sin contingencia} & \textbf{Con contingencia} \\
\midrule
Personal & 1.136~horas de cuatro integrantes. & 10.780.000 & 11.876.000 \\
Subcontrataciones & 26~horas de intérpretes de LSA a \$7.000 por hora; 32~horas de personal de salud sin costo. & 182.000 & 191.100 \\
Costes varios & Cuatro notebooks, internet, electricidad, cámara y micrófono. & 784.560 & 806.582,88 \\
Otros gastos & Material y equipo (viáticos de reunión), viajes y cuentas compartidas de un asistente de IA. & 217.200 & 221.544 \\
\midrule
\textbf{Total de costes y gastos} & Contingencia total: \$1.132.240. & \textbf{11.963.760} & \textbf{13.096.000} \\
Margen fijo (15~\%) & Sobre el total con contingencia. & & 1.964.400 \\
\textbf{Precio de venta (sin IVA)} & & & \textbf{15.060.400} \\
Precio de venta con IVA (21~\%) & & & 18.223.084 \\
\bottomrule
\end{tabularx}
\end{table}

El costo de personal se calculó con tarifas de \$10.000 por hora para las categorías de Líder Técnico y Desarrollador Backend, y de \$9.000 por hora para Desarrollador Frontend y para QA y Desarrollador Fullstack; las horas por categoría son 258 de Líder Técnico (132 y 126), 298 de Desarrollador Backend (166 y 132), 342 de Desarrollador Frontend (182 y 160) y 238 de QA y Fullstack (126 y 112). Las tasas de contingencia varían por rubro entre el 1~\% y el 11~\%.

Según el Plan Financiero, los ingresos provendrían de fondos de investigación y desarrollo (FONCyT y PROINE) y de la participación en el Premio INVAP a las mejores tesis de ingeniería, distribuidos en ocho pagos escalonados a lo largo de la ejecución: \$1.000.000, \$4.000.000, \$3.000.000, \$2.000.000, \$1.000.000, \$1.500.000, \$1.000.000 y \$2.000.000, que suman \$15.500.000. El Plan de Dirección agrega los aportes del cliente principal, el Sanatorio 9 de Julio, dentro de un esquema financiero mixto.

\chapter{Desarrollo del tema}

\section{Modelo de negocio}
\pendiente{propuesta de valor, cliente y sostenibilidad.}

\section{Desarrollo del producto}

\subsection{Análisis}
\subsubsection{Especificación de Requerimientos de Software}
\pendiente{fuente: TP7 (ERS IEEE 830) y Matriz de Trazabilidad.}
\subsubsection{Modelo de Casos de Uso}
\pendiente{fuente: TP7 (diagrama y modelos de casos de uso).}
\subsubsection{Modelo de Datos}
\pendiente{fuente: TP7 (modelo de dominio y modelo de datos).}

\subsection{Diseño}
\subsubsection{Documento de Arquitectura del Software}
\pendiente{vistas lógica y física, tipo de arquitectura, beneficios y consideraciones. Fuente: TP7.}
\subsubsection{Mockups}
\pendiente{fuente: TP7.}
\subsubsection{Diagrama Entidad Relación}
\pendiente{fuente: TP7.}
\subsubsection{Diagramas de Colaboración}
\pendiente{fuente: TP7.}

\subsection{Implementación}
\subsubsection{Tecnologías de implementación}
\pendiente{herramientas elegidas y comparativas: lenguajes, frameworks, modelo de IA, base de datos.}
\subsubsection{Modelo de implementación}
\pendiente{diagrama de despliegue, especificaciones técnicas y estrategia de implementación (Docker Compose).}

\subsection{Pruebas}
\subsubsection{Modelo de pruebas}
\pendiente{fuente: TP8 (Plan de Testing).}
\subsubsection{Plan de pruebas}
\pendiente{casos de prueba en ficha y resultados de los ciclos 1 y 2. Fuente: TP8.}

\chapter{Manual del usuario}\label{cap:manual-usuario}

Este manual está dirigido a las personas que usan LSA Salud en la atención: el paciente que se comunica en Lengua de Señas Argentina (LSA) y el personal de recepción o enfermería que lo atiende. No requiere conocimientos técnicos. La instalación y el mantenimiento del equipo se describen en el capítulo~\ref{cap:manual-instalacion}.

\section{Qué es LSA Salud}

LSA Salud es un traductor de LSA a texto y voz para la atención en guardia. El paciente hace señas frente a una cámara; el sistema las reconoce, arma una oración en español, la muestra en pantalla y la dice en voz alta. El personal responde escribiendo o dictando, y el paciente lee la respuesta en su pantalla. Todo el procesamiento ocurre en la propia computadora, sin enviar datos a servicios en la nube.

El sistema tiene dos pantallas, o estaciones, que se abren en el navegador y se sincronizan entre sí:
\begin{itemize}
  \item la \textbf{Estación del Paciente}, con la cámara, la conversación y las señas reconocidas;
  \item la \textbf{Estación de Recepción}, con la conversación, el cuadro para responder y la imagen de la cámara del paciente.
\end{itemize}

En esta versión el sistema reconoce cuatro señas: \emph{yo}, \emph{cabeza}, \emph{dolor} y \emph{respirar} (tabla~\ref{tab:especificaciones}). Combinadas, permiten expresar frases como ``yo dolor cabeza'', que el sistema traduce como ``Me duele la cabeza.''. Una seña fuera de ese vocabulario no se reconoce: la traducción es una ayuda para iniciar la comunicación y no reemplaza al intérprete humano. Ante cualquier duda sobre lo que el paciente quiso decir, hay que confirmarlo por otros medios.

\paragraph{Sobre las capturas de este manual.} Las imágenes se tomaron con la aplicación en funcionamiento, usando la cámara virtual de prueba del navegador, que muestra un patrón verde en lugar de una persona. Como esa cámara no muestra manos, el mensaje del paciente de los ejemplos (``Me duele la cabeza.'') es la oración que el sistema devuelve para la secuencia \emph{yo, dolor, cabeza} y se incorporó a la conversación sin hacer las señas. El resto de los elementos (consentimiento, botones, avisos y respuestas) son los que ve el usuario real.

\section{Inicio}

Al abrir la dirección del sistema en Google Chrome aparece la pantalla de inicio (figura~\ref{fig:man-inicio}) con el acceso a las estaciones. Normalmente el personal técnico deja ambas estaciones abiertas, una en cada monitor; en ese caso no es necesario usar esta pantalla. Quien necesite abrirlas debe usar el mismo navegador para las dos, porque la sincronización de la conversación solo funciona entre pestañas o ventanas del mismo navegador y perfil.

\figuraTP{manual/01_inicio}{0.9\textwidth}{Pantalla de inicio con los accesos a las estaciones}{fig:man-inicio}

La tarjeta \textbf{Captura de señas} corresponde al Modo Captura, una herramienta interna del equipo técnico para grabar muestras y reentrenar el modelo; no forma parte del uso en la atención (véase el capítulo~\ref{cap:manual-instalacion}).

\section{Consentimiento de uso de la cámara}

La primera vez que se abre una estación, el sistema muestra el diálogo \textbf{Consentimiento para usar la cámara} (figura~\ref{fig:man-consentimiento}). El texto explica que la cámara se usa solo para reconocer las señas y traducirlas a texto y voz, y que no se guarda ningún video ni imagen. Hay que leerlo con el paciente antes de continuar.

\figuraTP{manual/02_consentimiento}{0.9\textwidth}{Diálogo de consentimiento para usar la cámara}{fig:man-consentimiento}

\begin{itemize}
  \item \textbf{Acepto y continúo}: el paciente autoriza el uso de la cámara. Inmediatamente después, Chrome puede pedir permiso para usarla; hay que elegir \textbf{Permitir}.
  \item \textbf{No acepto}: la cámara no se activa y el traductor no funciona. El diálogo muestra el aviso \emph{``El traductor necesita la cámara para funcionar. Podés reconsiderar cuando quieras.''} y el botón \textbf{Reconsiderar} (figura~\ref{fig:man-rechazo}), que vuelve a mostrar la pregunta.
\end{itemize}

\figuraTP{manual/03_consentimiento_rechazado}{0.9\textwidth}{Diálogo tras elegir ``No acepto''}{fig:man-rechazo}

La decisión queda recordada en ese navegador, por lo que no se vuelve a preguntar en cada atención, y se registra como un evento de auditoría sin datos personales. Para cambiar una decisión ya guardada, el personal técnico debe borrar los datos del sitio en el navegador.

\section{Estación del Paciente}

\subsection{Qué se ve en pantalla}

La figura~\ref{fig:man-paciente} muestra la Estación del Paciente con una conversación en curso. Sus elementos son:
\begin{description}[style=nextline, leftmargin=1.5em]
  \item[Conversación (izquierda)] Muestra los mensajes. El título dice \emph{``Conversación''} y, cuando recepción cargó su nombre, \emph{``[nombre] te está atendiendo''}. Los mensajes del paciente aparecen a la derecha, en azul; los del personal, a la izquierda, en blanco. Si todavía no hay mensajes se lee \emph{``La conversación aparecerá aquí''}.
  \item[Captando señas (derecha)] Muestra la imagen de la cámara. El indicador \emph{``EN VIVO · Visión por computadora''} confirma que la cámara está activa. Mientras se activa, se lee \emph{``Activando cámara...''}.
  \item[Mostrar esqueleto] Interruptor que dibuja o quita las líneas que marcan manos, cuerpo y rostro sobre la imagen.
  \item[Mostrar palabras] Interruptor que muestra o esconde, sobre la imagen, las palabras que el sistema va reconociendo.
  \item[Top-3 Predicciones] Las tres señas que el sistema considera más probables en este momento, con su porcentaje de confianza. Es una ayuda para el equipo técnico; el paciente no necesita mirarla.
  \item[Limpiar Historial] Borra todos los mensajes de la conversación en ambas estaciones.
\end{description}

\figuraTP{manual/06_paciente_conversacion}{\textwidth}{Estación del Paciente con una conversación en curso}{fig:man-paciente}

\subsection{Cómo comunicarse}

\begin{enumerate}
  \item El paciente se ubica frente a la cámara, de modo que se vean el rostro, el torso y ambas manos. Conviene una iluminación pareja y evitar fondos con mucho movimiento.
  \item Hace las señas una a continuación de la otra, por ejemplo \emph{yo}, \emph{dolor}, \emph{cabeza}. Las palabras reconocidas se muestran sobre la imagen mientras \emph{Mostrar palabras} está activado.
  \item Al terminar, \textbf{baja las manos} y las deja en reposo frente al cuerpo. Ese reposo le indica al sistema que la oración terminó.
  \item Aparece el aviso \emph{``Generando respuesta de IA...''}. Pasados unos instantes, la oración en español se agrega a la conversación como mensaje del paciente.
  \item Un momento después de mostrarse el texto, el sistema lo dice en voz alta, de modo que el personal lo escuche. Si el sonido no se escucha, el texto igualmente queda en pantalla.
  \item El personal responde por escrito o por voz y el paciente lee la respuesta en la conversación. Las respuestas del personal no se leen en voz alta.
\end{enumerate}

Para hacer una nueva oración, el paciente repite el procedimiento. Si baja las manos sin haber hecho ninguna seña, el sistema descarta ese silencio y no genera ningún mensaje.

\subsection{Si la traducción es incorrecta}

Debajo de cada mensaje del paciente aparece el botón \textbf{Reportar traducción}. Si la oración no es lo que el paciente quiso decir, hay que tocarlo (figura~\ref{fig:man-reporte}): el mensaje queda tachado, marcado con \emph{``Traducción reportada''}, y se informa \emph{``Traducción reportada. Repetí la seña para enviar el mensaje de nuevo.''}. El paciente debe entonces repetir las señas. En la Estación de Recepción el mismo mensaje se marca como reportado y se avisa que no debe tomarse como válido.

\figuraTP{manual/09_paciente_reporte}{\textwidth}{Un mensaje reportado como incorrecto en la Estación del Paciente}{fig:man-reporte}

El reporte guarda el texto del mensaje, las palabras reconocidas y la estación desde la que se reportó, para que el equipo técnico mejore el reconocimiento. Si el reporte no se puede enviar, aparece \emph{``No se pudo enviar el reporte. Probá de nuevo.''}; puede intentarse otra vez, y si persiste, avisar al personal técnico.

\section{Estación de Recepción}

\subsection{Qué se ve en pantalla}

La figura~\ref{fig:man-recepcion} muestra la Estación de Recepción. Sus elementos son:
\begin{description}[style=nextline, leftmargin=1.5em]
  \item[Conversación con paciente] Lista de mensajes: los del paciente a la izquierda, en blanco, y las respuestas del personal a la derecha, en azul. Sin mensajes se lee \emph{``No hay mensajes en la conversación.''}.
  \item[Tu nombre] Nombre de quien atiende. Se cambia con el ícono del lápiz (\emph{Editar tu nombre}); se confirma con Enter o al salir del campo. El paciente lo ve en su pantalla como \emph{``[nombre] te está atendiendo''}. Queda guardado en el navegador.
  \item[Finalizar atención] Cierra la atención: borra la conversación en ambas estaciones. Solo actúa si hay mensajes. Debe usarse al terminar con cada paciente para que el siguiente no vea la conversación anterior.
  \item[Cuadro de respuesta] Campo \emph{``Escribir...''}, botón de micrófono y botón de envío.
  \item[Cámara del paciente] Imagen de la cámara del paciente, con los interruptores \emph{Mostrar esqueleto} y \emph{Mostrar palabras}. Mientras no llega la imagen se lee \emph{``Esperando la cámara del paciente''}.
\end{description}

\figuraTP{manual/08_recepcion_conversacion}{\textwidth}{Estación de Recepción con una conversación en curso}{fig:man-recepcion}

\subsection{Cómo atender}

\begin{enumerate}
  \item Antes de recibir al paciente, tocar el lápiz junto a \emph{Tu nombre}, escribir el nombre y confirmar con Enter.
  \item Leer los mensajes del paciente a medida que aparecen; también se escuchan en voz alta. El texto se acompaña con la imagen de la cámara para verificar lo que el paciente expresó con las señas.
  \item Responder de una de estas formas:
  \begin{itemize}
    \item \textbf{Por escrito:} escribir en \emph{Escribir...} y tocar el botón \textbf{Enviar} o presionar Enter. El botón está deshabilitado mientras el campo está vacío.
    \item \textbf{Por voz:} tocar el botón del micrófono (\emph{Dictar respuesta por voz}), hablar y volver a tocarlo (\emph{Detener dictado}) cuando se termina. Mientras se escucha, el botón se pone rojo y el campo no se puede editar. El texto dictado aparece en el campo y puede corregirse antes de enviarlo. La primera vez, Chrome solicita permiso para usar el micrófono.
  \end{itemize}
  \item El mensaje aparece de inmediato en la pantalla del paciente.
  \item Si una respuesta propia fue enviada por error, el ícono del cesto de basura (\emph{Eliminar mensaje}) la borra de ambas pantallas.
  \item Si el paciente reporta una traducción, o si el propio personal considera que es incorrecta, puede usar \textbf{Reportar traducción} en ese mensaje; el paciente repetirá la seña.
  \item Al terminar la atención, tocar \textbf{Finalizar atención}.
\end{enumerate}

\section{Mensajes de error y situaciones frecuentes}

La tabla~\ref{tab:man-errores} resume lo que puede ocurrir durante el uso y cómo actuar. Si el problema persiste, hay que avisar al personal técnico, que dispone de la guía de solución de problemas del capítulo~\ref{cap:manual-instalacion}.

\begin{table}[H]
  \centering
  \caption{Situaciones frecuentes durante el uso}\label{tab:man-errores}
  \small
  \begin{tabularx}{\textwidth}{>{\raggedright\arraybackslash}p{4.4cm} X}
    \toprule
    \textbf{Qué se ve} & \textbf{Qué hacer} \\
    \midrule
    \emph{No se pudo acceder a la cámara.} (figura~\ref{fig:man-camara}) & El navegador no tiene permiso para usar la cámara o la está usando otra aplicación. Permitir la cámara desde el ícono del candado junto a la dirección, cerrar las otras aplicaciones que la usen y recargar la página. \\
    \emph{Activando cámara...} que no termina & Esperar unos segundos; si no cambia, recargar la página. \\
    El diálogo de consentimiento con el botón \emph{Reconsiderar} & El paciente no aceptó el uso de la cámara. Explicarle el propósito y, si acepta, tocar \textbf{Reconsiderar} y luego \textbf{Acepto y continúo}. \\
    \emph{Esperando la cámara del paciente} en Recepción & La Estación del Paciente no está abierta o no tiene cámara. Abrirla en el mismo navegador. \\
    Las señas no producen ningún mensaje & Verificar que se vean las manos y el torso, que haya buena iluminación y que el paciente baje las manos al terminar la oración. Si el indicador \emph{Top-3 Predicciones} no cambia, avisar al personal técnico. \\
    La oración aparece pero no se escucha & Revisar el volumen y los altavoces. El texto igualmente queda en la conversación. \\
    \emph{El navegador no soporta reconocimiento de voz.} & El dictado solo funciona en Google Chrome. Responder por escrito o abrir la estación en Chrome. \\
    Los mensajes no llegan a la otra pantalla & Las dos estaciones deben estar abiertas en el mismo navegador y perfil, en la misma computadora. \\
    \bottomrule
  \end{tabularx}
\end{table}

\figuraTP{manual/04_camara_denegada}{\textwidth}{Estación del Paciente cuando no se pudo acceder a la cámara}{fig:man-camara}

\section{Privacidad}

LSA Salud se diseñó para proteger los datos del paciente:
\begin{itemize}
  \item \textbf{No se guardan imágenes ni video.} El navegador procesa la imagen de la cámara en el propio equipo y solo envía al sistema las coordenadas de los puntos de manos, cuerpo y rostro, que no permiten reconstruir la imagen.
  \item \textbf{No hay servicios en la nube.} El reconocimiento, la redacción de la oración y la voz se ejecutan en la misma computadora; el sistema funciona sin conexión a Internet.
  \item \textbf{La conversación no se guarda.} Los mensajes viven solo en la memoria del navegador y desaparecen al tocar \textbf{Finalizar atención} o \textbf{Limpiar Historial}, o al cerrar la pestaña.
  \item \textbf{Solo se conserva lo que se reporta.} Cuando alguien reporta una traducción, se guarda el texto de ese mensaje, las palabras reconocidas y la estación; no se guardan datos de identidad. El sistema también registra eventos de auditoría (por ejemplo, la aceptación del consentimiento y tiempos de respuesta), sin datos personales.
  \item \textbf{La cámara solo se activa con consentimiento.} Sin la aceptación del paciente, la cámara permanece apagada.
\end{itemize}

\chapter{Manual de instalación}
\pendiente{configuración de servidor y clientes: requisitos, Docker Compose y variables de entorno. Fuente: README del repositorio.}

\chapter{Incumbencias del Ingeniero en Sistemas de Información}

La tabla~\ref{tab:incumbencias} relaciona las quince incumbencias del título con lo realizado en el proyecto. Cuando el proyecto no ejercitó una incumbencia, se indica con «No aplica» y el motivo; cuando la ejercitó solo en parte, se aclara el alcance.

{\small
\begin{longtable}{>{\raggedright\arraybackslash}p{0.46\textwidth} >{\raggedright\arraybackslash}p{0.46\textwidth}}
    \caption{Incumbencias aplicadas en el proyecto}\label{tab:incumbencias} \\
    \toprule
    \textbf{Incumbencia} & \textbf{Justificación en el proyecto} \\
    \midrule
    \endfirsthead
    \toprule
    \textbf{Incumbencia} & \textbf{Justificación en el proyecto} \\
    \midrule
    \endhead
    1. Participar en la toma de decisiones estratégicas de una organización y asesorar, en concordancia con las mismas, acerca de las políticas de desarrollo de sistemas de información & Parcial. El equipo definió con el cliente el problema y los objetivos del proyecto (tabla~\ref{tab:objetivos}) y propuso decisiones de política tecnológica, como el procesamiento local sin servicios en la nube (tabla~\ref{tab:decisiones}) y el modelo de implementación y soporte anual (sección~\ref{sec:modelo-negocio}). \\ \midrule
    2. Evaluar, clasificar y seleccionar proyectos de sistemas de información y evaluar y seleccionar alternativas de asistencia interna & Parcial. El equipo comparó las soluciones existentes para la comunicación con personas sordas (sección~\ref{sec:similares}) y las metodologías de desarrollo aplicables (tabla~\ref{tab:metodologias}), y eligió su proyecto entre propuestas de tema (tabla~\ref{tab:hitos}). No clasificó una cartera de proyectos de una organización. \\ \midrule
    3. Planificar, efectuar y evaluar los estudios de factibilidad inherentes a todo proyecto de diseño de sistemas de información y de modificación o reemplazo de los mismos, así como los sistemas de computación asociados & Se estimó la factibilidad económica con el Plan Financiero (tabla~\ref{tab:presupuesto}), la factibilidad temporal con el cronograma (tabla~\ref{tab:hitos}) y los riesgos del proyecto (tabla~\ref{tab:riesgos}). La factibilidad técnica se verificó con el PMV, que cumplió la latencia y la exactitud previstas (tabla~\ref{tab:crit-res}). \\ \midrule
    4. Planificar, dirigir, ejecutar y controlar el relevamiento, análisis, diseño, desarrollo, implementación y prueba de sistemas de información & El proyecto recorrió todo el ciclo: relevamiento y análisis con la ERS y los casos de uso, diseño de arquitectura y de datos, implementación con Docker y pruebas de sistema en dos ciclos (capítulo~\ref{cap:desarrollo}), planificadas y controladas con el plan de gestión (capítulo~\ref{cap:gestion}). \\ \midrule
    5. Evaluar y seleccionar los sistemas de programación disponibles con miras a su utilización en sistemas de información & Se compararon y seleccionaron los lenguajes, bibliotecas y motores del sistema (MediaPipe, PyTorch, FastAPI, Ollama y Piper, entre otros) según los criterios de ejecución local y privacidad (tablas~\ref{tab:tecnologias} y~\ref{tab:comparativa}). \\ \midrule
    6. Evaluar y seleccionar, desde el punto de vista de los sistemas de información, los equipos de procesamiento y comunicación y los sistemas de base & Se definió que la estación opere con una PC, una cámara web y pantallas comerciales, sin hardware especializado (CRIT-06, tabla~\ref{tab:crit}), y se verificó el sistema sobre una PC comercial (tablas~\ref{tab:req-oficiales} y~\ref{tab:especificaciones}). Los riesgos de hardware se trataron en la tabla~\ref{tab:respuesta}. \\ \midrule
    7. Organizar y dirigir el área de sistemas; determinar el perfil de los recursos humanos necesarios y contribuir a su selección y formación & No aplica: el proyecto no incluyó la dirección de un área de sistemas de una organización ni la selección de personal. Solo se organizó el propio equipo, con roles y esfuerzo por paquete de trabajo (tablas~\ref{tab:roles} y~\ref{tab:edt}). \\ \midrule
    8. Participar en la elaboración de programas de capacitación para la utilización de sistemas de información & Se elaboraron el manual de usuario y el de instalación (capítulos~\ref{cap:manual-usuario} y~\ref{cap:manual-instalacion}), y la capacitación del personal se incluyó en el componente de implementación del modelo de ingresos (sección~\ref{sec:modelo-negocio}). Todavía no se dictó capacitación en el sanatorio. \\ \midrule
    9. Determinar y controlar el cumplimiento de las pautas técnicas que rigen el funcionamiento y la utilización de recursos informáticos en cada organización & No aplica en esta etapa: el PMV no se implantó en la infraestructura del cliente, por lo que no se controlaron pautas técnicas de la organización. El manual de instalación (capítulo~\ref{cap:manual-instalacion}) documenta los requisitos de operación del sistema. \\ \midrule
    10. Elaborar métodos y normas a seguir en cuestiones de seguridad y privacidad de la información procesada y la generada por los sistemas de información; participar en la determinación de las acciones a seguir en esta materia y evaluar su aplicación & La privacidad se incorporó desde el diseño: procesamiento local, consentimiento explícito previo a la cámara (CRIT-07) y ausencia de almacenamiento persistente de imágenes (CRIT-04). Su aplicación se evaluó con los casos de prueba de consentimiento y de privacidad y seguridad (tablas~\ref{tab:crit-res} y~\ref{tab:casos}). \\ \midrule
    11. Elaborar métodos y normas a seguir en cuestión de salvaguardia y control de los recursos físicos y lógicos de un sistema de computación; participar en la determinación de las acciones a seguir en esta materia y evaluar su aplicación & Parcial, solo sobre recursos lógicos: se restringieron los orígenes permitidos a la red local, se validaron las entradas y se limitó la información persistida a los reportes y los eventos de auditoría anónimos (tablas~\ref{tab:decisiones} y~\ref{tab:casos}). No se trataron recursos físicos de una organización. \\ \midrule
    12. Desarrollar modelos de simulación, sistemas expertos y otros sistemas informáticos destinados a la resolución de problemas y asesorar en su aplicación & Se desarrolló un sistema de reconocimiento de señas basado en una red neuronal LSTM, con reconstrucción de oraciones y síntesis de voz locales (tabla~\ref{tab:tecnologias}), evaluado en la sección~\ref{sec:res-pruebas}, y se elaboró el manual para su aplicación (capítulo~\ref{cap:manual-usuario}). \\ \midrule
    13. Realizar auditorías en áreas de sistemas y centros de cómputos, así como en los sistemas de información utilizados & No aplica: el proyecto no incluyó auditorías de un área de sistemas ni de sistemas ajenos. Las pruebas de sistema (sección~\ref{sec:modelo-pruebas}) verifican el producto propio, y el registro de eventos anónimos (RF-04) habilita auditorías futuras. \\ \midrule
    14. Realizar arbitrajes, pericias y tasaciones referidas a los sistemas de información y a los medios de procesamiento de datos & No aplica: el proyecto no incluyó arbitrajes, pericias ni tasaciones. El precio de venta estimado (tabla~\ref{tab:presupuesto}) corresponde al propio proyecto y no a una tasación de sistemas de terceros. \\ \midrule
    15. Realizar estudios e investigaciones conducentes a la creación y mejoramiento de técnicas de desarrollo de sistemas de información y nuevas aplicaciones de la tecnología informática existente & Se realizó una investigación del estado del arte (capítulo~\ref{cap:estado-arte}) y se aplicó tecnología existente (visión por computadora y redes de secuencias) a una nueva aplicación: la traducción de un vocabulario clínico de la LSA, con un dataset propio y un modo de captura para ampliarlo. \\
    \bottomrule
\end{longtable}
}

\chapter{Resultados obtenidos}\label{cap:resultados}

Este capítulo contrasta lo planificado con lo obtenido: los objetivos específicos del capítulo de Objetivos (tabla~\ref{tab:objetivos}), los criterios de aceptación del capítulo de Gestión (tabla~\ref{tab:crit}), el producto entregado, las métricas de las pruebas y el desempeño de la gestión del proyecto. Todos los valores provienen de las pruebas del Producto Mínimo Viable (PMV) descriptas en la sección~\ref{sec:pruebas-plan} y del seguimiento del Informe Final. Los objetivos que requieren operación sostenida en la guardia del sanatorio, etapa posterior al PMV, se marcan como no evaluables y no se estiman a partir de otros valores.

\section{Cumplimiento de los objetivos específicos}\label{sec:res-objetivos}

La tabla~\ref{tab:res-objetivos} contrasta cada objetivo específico con lo obtenido en el PMV. El estado \emph{Cumplido} indica que la meta se midió y se alcanzó en las condiciones descriptas; \emph{Cumplido parcialmente}, que solo una parte de la meta pudo verificarse; y \emph{No evaluable en el PMV}, que la medición exige una condición que el PMV no tuvo, indicada en cada caso.

{\footnotesize
\setlength{\tabcolsep}{3pt}
\begin{longtable}{@{}p{0.9cm}p{3.3cm}p{5.0cm}p{2.7cm}p{3.1cm}@{}}
\caption{Contraste entre los objetivos específicos y los resultados del PMV.}\label{tab:res-objetivos}\\
\toprule
\textbf{ID} & \textbf{Meta} & \textbf{Resultado obtenido en el PMV} & \textbf{Evidencia} & \textbf{Estado} \\
\midrule
\endfirsthead
\toprule
\textbf{ID} & \textbf{Meta} & \textbf{Resultado obtenido en el PMV} & \textbf{Evidencia} & \textbf{Estado} \\
\midrule
\endhead
\midrule
\multicolumn{5}{r}{\emph{Continúa en la página siguiente}}\\
\endfoot
\bottomrule
\endlastfoot
OE01 & Reducir al menos un 80~\% el tiempo de las interacciones administrativas y de triaje (de 4~minutos a menos de 1~minuto). & No se midió el tiempo de una interacción completa con pacientes reales. Solo se midió la latencia del sistema (p50 de 166~ms y p95 de 945~ms), que es una parte de ese tiempo. & Tabla~\ref{tab:metricas}; sección~\ref{sec:pruebas-plan} & No evaluable en el PMV: requiere registros de uso en la guardia. \\
OE02 & Contribuir a reducir al menos un 50~\% la brecha de permanencia entre pacientes sordos y oyentes. & No se midió la permanencia en urgencias. & Tabla~\ref{tab:hitos} (H09 y H10) & No evaluable en el PMV: requiere el piloto en guardia y un año de uso operativo. \\
OE03 & Recuperar un mínimo estimado de 150~horas anuales del personal de guardia. & No se registraron tiempos de atención que permitan la extrapolación. & Tabla~\ref{tab:hitos} (H09 y H10) & No evaluable en el PMV: requiere operación sostenida en la guardia. \\
OE04 & Elevar la comprensión clínica de un valor de referencia cercano al 18~\% a más del 90~\%. & No se aplicaron encuestas ni preguntas de control a pacientes. Las pruebas con una persona real verificaron la calidad de la traducción, no la comprensión de diagnósticos. & Sección~\ref{sec:pruebas-plan} & No evaluable en el PMV: requiere encuestas posatención con pacientes usuarios de LSA. \\
OE05 & Reducir la readmisión a 30~días de cerca del 24,3~\% a menos del 15~\%. & No se realizó análisis longitudinal. & Tabla~\ref{tab:hitos} (H10) & No evaluable en el PMV: requiere 24~meses de funcionamiento. \\
OE06 & Precisión de al menos 90~\% en contextos clínicos controlados y 85~\% en entornos reales de guardia. & Exactitud media del 98,67~\% en validación cruzada de cinco particiones (96,67~\% a 100~\% por partición) y \emph{recall} por clase de 96,7~\% a 100~\%. Las particiones provienen de una sola sesión de grabación, por lo que el valor puede sobrestimar la generalización. No se midió la precisión en la guardia ni el \emph{F1-score}. & Tabla~\ref{tab:metricas}; tabla~\ref{tab:crit-res} (CRIT-02 y CRIT-03) & Cumplido parcialmente: cumplido en condición controlada; el entorno real no se evaluó. \\
OE07 & Tiempo de respuesta inferior a 1~segundo en la interpretación en tiempo real. & Latencia percibida de 166~ms (p50) y 945~ms (p95) sobre 55~oraciones, medida en uso real con una persona. El margen del p95 es estrecho. & Tabla~\ref{tab:metricas}; tabla~\ref{tab:rendimiento} & Cumplido, con la medición a repetir en la estación definitiva. \\
OE08 & Que al menos el 70~\% del personal de recepción y triaje use el sistema. & No hubo personal usuario en operación. & Tabla~\ref{tab:hitos} (H09 y H10) & No evaluable en el PMV: requiere registros de uso durante 6~meses de implementación. \\
OE09 & Satisfacción superior al 85~\% en pacientes usuarios de LSA. & No se aplicaron encuestas de satisfacción. & Tabla~\ref{tab:hitos} (H09 y H10) & No evaluable en el PMV: requiere encuestas posconsulta con pacientes. \\
OE10 & 100~\% de las interacciones sin almacenamiento persistente de imágenes ni datos sensibles. & Por diseño, al backend solo llegan coordenadas de landmarks. La prueba TC-LSA-013-01 verificó que las ventanas procesadas no modificaron ningún archivo y que solo se persiste el texto del mensaje que un usuario reporta, además de eventos de auditoría anónimos. & Tablas~\ref{tab:casos}, \ref{tab:req-oficiales} (RNF-03) y~\ref{tab:comparativa} & Cumplido parcialmente: verificado en el PMV por pruebas propias; la verificación a lo largo del ciclo de vida y la revisión normativa son posteriores. \\
\end{longtable}
}

En síntesis, de los diez objetivos específicos, uno (OE07) se cumplió, dos (OE06 y OE10) se cumplieron parcialmente y siete (OE01 a OE05, OE08 y OE09) no pueden evaluarse con el PMV. Estos últimos son los objetivos de impacto operativo y clínico, cuya medición depende del piloto en la guardia del Sanatorio 9 de Julio, que el Informe Final registra como etapa posterior al PMV.

\section{Cumplimiento de los criterios de aceptación}\label{sec:res-criterios}

La tabla~\ref{tab:crit-res} contrasta los siete criterios de aceptación (tabla~\ref{tab:crit}) con los resultados del PMV.

{\footnotesize
\setlength{\tabcolsep}{3pt}
\begin{longtable}{@{}p{1.4cm}p{6.0cm}p{3.6cm}p{4.0cm}@{}}
\caption{Contraste entre los criterios de aceptación y los resultados del PMV.}\label{tab:crit-res}\\
\toprule
\textbf{ID} & \textbf{Resultado obtenido} & \textbf{Evidencia} & \textbf{Estado} \\
\midrule
\endfirsthead
\toprule
\textbf{ID} & \textbf{Resultado obtenido} & \textbf{Evidencia} & \textbf{Estado} \\
\midrule
\endhead
\midrule
\multicolumn{4}{r}{\emph{Continúa en la página siguiente}}\\
\endfoot
\bottomrule
\endlastfoot
CRIT-01 & Latencia percibida de 945~ms (p95) y 166~ms (p50) sobre 55~oraciones en uso real, inferior a 1~segundo. & Tablas~\ref{tab:metricas} y~\ref{tab:rendimiento} & Cumplido. \\
CRIT-02 & Exactitud media del 98,67~\% en validación cruzada de cinco particiones, con \emph{recall} por clase de 96,7~\% a 100~\%. No se informó el \emph{F1-score}. La condición controlada es una única sesión de grabación. & Tabla~\ref{tab:metricas}; tabla~\ref{tab:casos} (TC-LSA-015-01) & Cumplido, con la salvedad de generalización. \\
CRIT-03 & No se realizaron pruebas operativas en la guardia real del sanatorio. & Tabla~\ref{tab:hitos} (H09) & No evaluable en el PMV: requiere el piloto en guardia. \\
CRIT-04 & Las ventanas procesadas en la prueba no modificaron ningún archivo; el backend solo recibe coordenadas y la única información persistida es el texto de los mensajes reportados y los eventos de auditoría anónimos. & Tablas~\ref{tab:casos} (TC-LSA-013-01), \ref{tab:req-oficiales} (RNF-03) y~\ref{tab:decisiones} & Cumplido en el PMV, verificado por la prueba de sistema y no por una auditoría independiente. \\
CRIT-05 & Las señas del PMV provienen de Spreadthesign y no cuentan con aval formal de intérpretes de LSA. & Tabla~\ref{tab:req-oficiales} (RNF-06) & No cumplido en el PMV: queda para la etapa de sistema completo. \\
CRIT-06 & La estación de prueba usó una PC y una cámara web comerciales, sin hardware médico. Las pantallas de doble cara no forman parte del PMV. & Tabla~\ref{tab:req-oficiales} (RNF-04 y RNF-07); tabla~\ref{tab:especificaciones} & Cumplido parcialmente: componentes estándar verificados; la estación física, pendiente. \\
CRIT-07 & El diálogo de consentimiento aparece antes de activar la cámara; sin aceptación, la cámara permanece apagada. & Tablas~\ref{tab:req-oficiales} (RF-03) y~\ref{tab:casos} (TC-LSA-009-01); capítulo~\ref{cap:manual-usuario} & Cumplido. \\
\end{longtable}
}

\section{Resultados del producto}\label{sec:res-producto}

El PMV entregado se ejecuta en una sola computadora con Docker Compose (capítulo~\ref{cap:manual-instalacion}) y comprende:
\begin{itemize}
  \item \textbf{Dos estaciones de cliente}, Paciente y Recepción. En la primera, la persona usuaria de LSA hace señas frente a la cámara, con consentimiento previo, y el sistema las traduce a texto y voz en español; en la segunda, el personal lee el mensaje y responde por texto o por dictado (capítulo~\ref{cap:manual-usuario}).
  \item \textbf{Una herramienta interna del equipo}, el Modo Captura, que permite grabar señas nuevas y reentrenar el modelo. No es una estación de cliente.
  \item \textbf{Un vocabulario de cuatro señas} (yo, cabeza, dolor y respirar) más la clase de reposo, tomadas de Spreadthesign y ampliables con el Modo Captura (tabla~\ref{tab:especificaciones}).
  \item \textbf{El reporte de traducciones incorrectas (RF-06)}, incorporado a pedido del cotutor, que desde ambas estaciones registra el mensaje reportado como insumo para reentrenar el modelo.
\end{itemize}

Según la tabla~\ref{tab:req-oficiales}, los seis requisitos funcionales (RF-01 a RF-06) están implementados y, de los siete no funcionales, RNF-01 a RNF-05 se cumplen en el PMV con el alcance indicado en esa tabla. RNF-06 (validación lingüística por intérpretes oficiales) y RNF-07 (estación fija de doble cara) corresponden a la etapa de sistema completo. Además, en el caso TC-LSA-015-01 el sistema siguió funcionando sin Internet, sin peticiones a hosts externos, y el análisis con axe-core no halló violaciones críticas ni graves de accesibilidad en las cuatro pantallas.

\section{Resultados de las pruebas}\label{sec:res-pruebas}

Las pruebas de sistema se ejecutaron en dos ciclos (20 y 21 de octubre de 2026). La tabla~\ref{tab:metricas} reúne las métricas principales; el detalle de los casos y de los defectos está en la sección~\ref{sec:pruebas-plan} y en el anexo~\ref{anx:defectos}.

\begin{table}[htbp]
\centering
\caption{Métricas de las pruebas del PMV.}\label{tab:metricas}
\footnotesize
\begin{tabularx}{\textwidth}{@{}p{4.3cm}p{3.2cm}p{2.6cm}X@{}}
\toprule
\textbf{Métrica} & \textbf{Resultado} & \textbf{Meta} & \textbf{Observación} \\
\midrule
Casos de prueba aprobados & Ciclo~1: 2 de 15; ciclo~2: 15 de 15 & 15 de 15 & Los 15 casos cubren los 12 casos de uso y tres grupos de requisitos no funcionales. \\
Defectos & 28 detectados (19 en el ciclo~1 y 9 en el ciclo~2); 28 cerrados & Ninguno abierto de severidad 1 o 2 & Tabla~\ref{tab:defectos-resumen}: 3 críticos, 15 altos, 8 medios y 2 bajos. \\
Exactitud en validación cruzada & 98,67~\% (96,67~\% a 100~\% por partición; desvío de 1,63) & Al menos 90~\% & Cinco particiones estratificadas, 150~muestras, dos errores en total. Una sola sesión de grabación. \\
Latencia percibida, p95 & 945~ms & Menos de 1~s & 55~oraciones medidas el 21/10/2026; el margen es estrecho. \\
Latencia percibida, p50 & 166~ms & 600~ms (ruta con plantilla) & La traducción con modelo de lenguaje tiene un p95 de 573~ms. \\
Inferencia del modelo, p95 & De 0,71~ms a 0,68~ms entre versiones & Regresión menor al 10~\% & Mismas 100~ventanas. \\
\bottomrule
\end{tabularx}
\end{table}

Estos resultados tienen tres límites que deben leerse junto con las cifras. Primero, la exactitud proviene de una única sesión de grabación (misma persona, iluminación y día), por lo que no mide la generalización a otras personas y condiciones; la medición con otras personas es una de las sugerencias del ciclo~2. Segundo, las comprobaciones manuales del ciclo~2 las realizó una persona, y la latencia debe repetirse en la estación definitiva. Tercero, el vocabulario es de cuatro señas, por lo que los resultados no se extrapolan a un vocabulario clínico más amplio.

\section{Resultados de la gestión del proyecto}\label{sec:res-gestion}

\paragraph{Cronograma.} Según la tabla~\ref{tab:hitos}, el balance acumulado de las fechas del Informe Final es de nueve días de adelanto respecto de las fechas originales. Cuatro hitos tuvieron demora positiva: el diseño de la arquitectura (21~días) y las pruebas piloto en entorno controlado (66~días), ambos por la reprogramación del cronograma aprobado, además del Plan de Riesgos (7~días) y la construcción del dataset (2~días). El PMV se desarrolló dentro de la fecha del Acta (H06, 13~días antes), y las pruebas formales se realizaron los días 20 y 21 de octubre. Los hitos de implementación piloto en guardia (H09), evaluación de impacto operativo (H10) y presentación de la versión final con manual de uso (H11) quedaron fuera del PMV.

\paragraph{Alcance y cambios.} El alcance del PMV se ajustó mediante decisiones registradas en el Informe Final: la ampliación de la EDT a 36 paquetes de trabajo por observación del cotutor, la incorporación del requisito RF-06, la construcción de un dataset propio con el Modo Captura en lugar de ampliar LSA64, el uso de un diccionario de referencia (Spreadthesign) en lugar del aval de intérpretes (RNF-06) y la ejecución local con Docker en lugar de un entorno en la nube. Los riesgos de privacidad, latencia y portabilidad se atendieron con medidas que se materializaron en el producto (capítulo~\ref{cap:gestion}).

\paragraph{Costos.} El Informe Final no registra costos reales de ejecución, por lo que no es posible contrastar el gasto con el presupuesto. Las cifras de la tabla~\ref{tab:presupuesto} son la estimación del Plan Financiero y no resultados.

\chapter{Conclusiones}\label{cap:conclusiones}

Este capítulo cierra el informe: valora en qué medida el proyecto cumplió su objetivo general, resume lo aprendido en la gestión, declara con franqueza las limitaciones del Producto Mínimo Viable (PMV) y propone el trabajo futuro. Las cifras provienen del capítulo de Resultados (capítulo~\ref{cap:resultados}) y no se reproducen más que lo necesario.

\section{Conclusiones generales}\label{sec:concl-generales}

El objetivo general (capítulo de Objetivos) consistía en desarrollar e implementar un sistema basado en inteligencia artificial y visión por computadora que facilite la comunicación entre personas usuarias de la Lengua de Señas Argentina (LSA) y el personal de salud. Con el alcance de un PMV, ese objetivo se cumplió \emph{en su parte de desarrollo}: existe un sistema funcional que traduce señas a texto y voz en español y permite al personal responder por texto o dictado. La parte de \emph{implementación} en la guardia del Sanatorio 9 de Julio no se realizó, por lo que el impacto sobre la accesibilidad, la calidad de la interacción y la equidad en el acceso, que son el fin último del objetivo general, no pudo medirse (sección~\ref{sec:res-objetivos}).

Lo que el proyecto sí demostró, con evidencia propia, es lo siguiente:
\begin{itemize}
  \item \textbf{Reconocimiento local y en tiempo real con hardware comercial.} Los 15 casos de prueba quedaron aprobados al cierre del ciclo~2 y la latencia percibida cumplió la meta de 1~segundo (OE07, CRIT-01), usando una PC y una cámara web comunes, sin hardware médico ni servicios en la nube (tabla~\ref{tab:metricas}). Esto es coherente con lo anticipado en el estado del arte sobre el reconocimiento de pose en el dispositivo (capítulo~\ref{cap:estado-arte}).
  \item \textbf{Privacidad por diseño.} Al backend solo llegan coordenadas de \emph{landmarks}; las imágenes no salen del navegador, la cámara se activa únicamente tras el consentimiento y el sistema funcionó sin conexión a Internet (CRIT-04 y CRIT-07, tabla~\ref{tab:crit-res}). La verificación fue propia y no una auditoría independiente.
  \item \textbf{Canal bidireccional.} El flujo completo, desde la seña hasta la respuesta del personal, está implementado en dos estaciones (Paciente y Recepción), con una herramienta interna, el Modo Captura, que permite ampliar el vocabulario y reentrenar el modelo (sección~\ref{sec:res-producto}).
  \item \textbf{Exactitud en condición controlada.} La validación cruzada arrojó una exactitud media del 98,67~\% sobre el vocabulario del PMV, por encima de la meta del 90~\% (OE06, CRIT-02), con la salvedad de generalización que se detalla más abajo.
\end{itemize}

En síntesis, de los diez objetivos específicos uno se cumplió, dos se cumplieron parcialmente y siete no son evaluables con un PMV (tabla~\ref{tab:res-objetivos}). Esa distribución no es un fracaso del proyecto sino la consecuencia de su alcance: los objetivos de impacto operativo y clínico exigen operación sostenida en la guardia, una etapa que el plan registró desde el inicio como posterior al PMV.

\section{Conclusiones sobre la gestión}\label{sec:concl-gestion}

El proyecto se gestionó con la Guía del PMBOK como marco de planificación y control y con desarrollo iterativo e incremental del producto (capítulo~\ref{cap:gestion}). Esa combinación permitió disponer de una versión mínima funcional temprana, que hizo posible medir latencia y precisión antes de las pruebas formales, y conservar la trazabilidad entre requisitos, casos de prueba y criterios de aceptación que exigían los entregables de la cátedra.

En términos de cronograma, el balance acumulado fue de nueve días de adelanto sobre las fechas originales del Informe Final, aunque con hitos demorados (sección~\ref{sec:res-gestion}); el PMV estuvo dentro de la fecha del Acta y las pruebas formales se ejecutaron en dos ciclos. El segundo ciclo cerró los defectos abiertos por el primero, lo que muestra que el proceso de prueba cumplió su función: el ciclo~1 aprobó 2 de 15 casos y detectó defectos críticos que el ciclo~2 permitió verificar como corregidos (anexo~\ref{anx:defectos}).

Funcionó bien: acotar el alcance a un vocabulario pequeño y decidirlo de forma explícita; atender temprano los riesgos de privacidad, latencia y portabilidad con medidas que terminaron materializadas en el producto; y construir un dataset propio con el Modo Captura en lugar de depender de datos ajenos. Quedó pendiente el registro de costos reales, que el Informe Final no incluye, por lo que no es posible contrastar el gasto con el presupuesto.

\section{Limitaciones del PMV}\label{sec:concl-limitaciones}

Las siguientes limitaciones deben leerse junto con cualquier cifra del informe.

\begin{itemize}
  \item \textbf{Vocabulario mínimo.} El sistema reconoce cuatro señas (yo, cabeza, dolor y respirar) más la clase de reposo. Los resultados no se extrapolan a un vocabulario clínico más amplio.
  \item \textbf{Una sola sesión de grabación.} Las particiones de la validación cruzada provienen de una única sesión (misma persona, iluminación y día). La exactitud del 98,67~\% puede por ello sobrestimar la generalización a otras personas y condiciones, y no se midió el \emph{F1-score}.
  \item \textbf{Margen de latencia estrecho.} El p95 de 945~ms cumple la meta de 1~s, pero con poca holgura, y se midió con una persona en la estación de prueba. Debe repetirse en la estación definitiva.
  \item \textbf{Sin validación lingüística.} Las señas provienen de un diccionario de referencia (Spreadthesign) y no cuentan con el aval de intérpretes de LSA ni de la comunidad sorda, por lo que CRIT-05 no se cumplió.
  \item \textbf{Sin piloto en la guardia.} Al no haberse operado en el sanatorio, OE01 a OE05, OE08 y OE09 no son evaluables, y no se evaluó la precisión en entorno real (parte de OE06). Las líneas de base de OE02, OE04 y OE05 provienen de estudios en otras poblaciones (sección~\ref{sec:guardias}); su primera medición local forma parte de esa etapa.
  \item \textbf{Estación física pendiente.} La estación fija de doble cara (RNF-07) no forma parte del PMV; la prueba usó componentes estándar (CRIT-06).
  \item \textbf{Verificación de privacidad propia.} OE10 se verificó mediante pruebas del equipo; una auditoría independiente y la revisión normativa son posteriores.
\end{itemize}

\section{Trabajo futuro}\label{sec:concl-futuro}

Se ordenan por prioridad, de lo que condiciona la validez de los resultados actuales a lo que amplía el alcance del producto.

\begin{enumerate}
  \item \textbf{Ampliar el vocabulario y la diversidad de las grabaciones.} Incorporar señas del triaje de baja complejidad mediante el Modo Captura, con grabaciones de varias personas, incluidas personas sordas usuarias de LSA, en distintas condiciones de iluminación y cámara. Es la medida que más directamente reduce el riesgo de sobreestimar la generalización.
  \item \textbf{Validar con intérpretes y con la comunidad.} Someter las señas y las traducciones a la revisión de intérpretes de LSA y de personas de la comunidad sorda, para cumplir CRIT-05 y reemplazar la fuente de referencia por un vocabulario avalado.
  \item \textbf{Realizar un piloto en la guardia del sanatorio.} Operar el sistema en el servicio, medir localmente las líneas de base y los indicadores de OE01 a OE05, OE08 y OE09, y evaluar la precisión en entorno real, con las salvaguardas éticas y de consentimiento correspondientes.
  \item \textbf{Repetir la medición de latencia en la estación definitiva} y con más de una persona, para confirmar que el margen del p95 se sostiene.
  \item \textbf{Construir la estación fija de doble cara} (sistema completo), con el hardware y la verificación que exige RNF-07.
  \item \textbf{Incorporar componentes no manuales.} Aprovechar expresión facial y posición de cabeza y torso, que el estado del arte señala como parte de la gramática de la LSA y que los sistemas suelen aprovechar poco.
  \item \textbf{Pasar del reconocimiento de señas aisladas al de oraciones continuas}, de modo que la traducción respete la gramática propia de la LSA y no solo el orden de las señas.
\end{enumerate}

\section{Cierre}\label{sec:concl-cierre}

La Ley 27.710 reconoce a la LSA como lengua natural y originaria, pero deja sin resolver cómo garantizar ese derecho en cada institución de salud. LSA Salud no resuelve ese problema por sí sola ni reemplaza a los intérpretes humanos; aporta una prueba de concepto de que un canal de comunicación local, que no almacena imágenes y corre en hardware común, es técnicamente viable. Que llegue a ser una herramienta útil para las personas sordas que se atienden en una guardia depende de lo que este informe deja explícitamente pendiente: la validación con la comunidad y la prueba en el servicio.


\printbibliography[heading=bibintoc, title={Bibliografía}]

\appendix
\chapter{Anexos}

\pendiente{enlace al repositorio del código y documentos de la cátedra; las planillas completas de casos de prueba (fichas de los 15 casos) y de riesgos no se reproducen en este informe.}

\section{Matriz de trazabilidad de requisitos}\label{anx:trazabilidad}

La tabla~\ref{tab:trazabilidad} vincula cada sub-requisito funcional de la ERS con los casos de uso que lo realizan y con los requisitos no funcionales relacionados. La tabla~\ref{tab:trazabilidad-rnf} indica los sub-requisitos técnicos de cada requisito no funcional de origen.

\begin{longtable}{>{\raggedright\arraybackslash}p{3.2cm} >{\raggedright\arraybackslash}p{5.0cm} >{\raggedright\arraybackslash}p{6.4cm}}
  \caption{Matriz de trazabilidad entre requisitos funcionales, casos de uso y requisitos no funcionales}\label{tab:trazabilidad}\\
  \toprule
  \textbf{Requisito funcional} & \textbf{Casos de uso} & \textbf{Requisitos no funcionales} \\
  \midrule
  \endfirsthead
  \toprule
  \textbf{Requisito funcional} & \textbf{Casos de uso} & \textbf{Requisitos no funcionales} \\
  \midrule
  \endhead
  RF-01.1 & CU-01 & RNF-T-01 \\
  RF-01.2 & CU-01 & RNF-T-01, RNF-01.1 \\
  RF-01.3 & CU-01, CU-02 & RNF-01.1 \\
  RF-01.4 & CU-01 & RNF-T-05, RNF-01.1 \\
  RF-01.5 & CU-02 & RNF-T-01, RNF-01.1 \\
  RF-01.6 & CU-03 & RNF-T-01, RNF-01.1 \\
  RF-01.7 & CU-01, CU-04 & RNF-T-05 \\
  RF-01.8 & CU-02 & --- \\
  RF-02.1 & CU-04 & --- \\
  RF-02.2 & CU-05 & RNF-T-04 (compatibilidad de navegador) \\
  RF-T-01 & CU-07 & --- \\
  RF-T-02 & CU-07 & --- \\
  RF-T-03 & CU-06 & --- \\
  RF-T-04 & CU-06 & RNF-03.1 \\
  RF-T-05 & CU-08 & --- \\
  RF-T-06 & Precondición de CU-01 y CU-04 & --- \\
  \midrule
  RF-03.1 & CU-09 & RNF-03.1 \\
  RF-04.1 & CU-11 & RNF-03.1 \\
  RF-06.1 & CU-10 & RNF-03.1 \\
  \bottomrule
\end{longtable}

\begin{table}[H]
  \centering
  \caption{Requisitos no funcionales de origen y sus sub-requisitos}\label{tab:trazabilidad-rnf}
  \small
  \begin{tabularx}{\textwidth}{>{\raggedright\arraybackslash}p{5.2cm} X}
    \toprule
    \textbf{Requisito no funcional} & \textbf{Sub-requisito técnico} \\
    \midrule
    RNF-01 (latencia) & RNF-01.1 \\
    RNF-02 (precisión) & RNF-02.1 \\
    RNF-03 (seguridad y privacidad) & RNF-03.1 \\
    RNF-04 (hardware) & RNF-04.1, compartido con RNF-07.1 \\
    RNF-05 (dataset especializado) & RNF-05.1 \\
    RNF-06 (fuente lingüística) & RNF-06.1 \\
    RNF-07 (estación fija) & RNF-07.1, compartido con RNF-04.1 \\
    \bottomrule
  \end{tabularx}
\end{table}

\section{Diccionario de datos de SQLite}\label{anx:diccionario}

Las tablas de la base de datos \texttt{data/lsa\_salud.db} implementan el modelo lógico de la figura~\ref{fig:modelo-logico}. La tabla~\ref{tab:dic-sqlite} detalla sus atributos y restricciones; las marcas de tiempo se expresan en formato ISO~8601 UTC.

\begin{longtable}{>{\raggedright\arraybackslash}p{3.8cm} >{\raggedright\arraybackslash}p{2.9cm} >{\raggedright\arraybackslash}p{7.9cm}}
  \caption{Diccionario de datos de las tablas de SQLite}\label{tab:dic-sqlite}\\
  \toprule
  \textbf{Atributo} & \textbf{Tipo} & \textbf{Restricciones y descripción} \\
  \midrule
  \endfirsthead
  \toprule
  \textbf{Atributo} & \textbf{Tipo} & \textbf{Restricciones y descripción} \\
  \midrule
  \endhead
  \multicolumn{3}{l}{\textbf{Tabla \texttt{reported\_message}}} \\
  \texttt{id} & \texttt{INTEGER} & Clave primaria autoincremental. \\
  \texttt{client\_message\_id} & \texttt{INTEGER} & Identificador asignado por el frontend; admite nulo. Junto con \texttt{text} es único, para no duplicar el mensaje cuando ambas caras lo reportan. \\
  \texttt{text} & \texttt{TEXT} & Texto del mensaje señalado, de 1 a 500 caracteres. \\
  \texttt{first\_reported\_at} & \texttt{TEXT} & Fecha y hora del primer reporte; obligatorio. \\
  \midrule
  \multicolumn{3}{l}{\textbf{Tabla \texttt{recognized\_word}}} \\
  \texttt{message\_id} & \texttt{INTEGER} & Clave foránea hacia \texttt{reported\_message}, con eliminación en cascada; forma parte de la clave primaria. \\
  \texttt{position} & \texttt{INTEGER} & Posición de la palabra en la frase, mayor o igual a cero; forma parte de la clave primaria. \\
  \texttt{word} & \texttt{TEXT} & Etiqueta de la seña reconocida, de 1 a 32 caracteres. \\
  \midrule
  \multicolumn{3}{l}{\textbf{Tabla \texttt{translation\_report}}} \\
  \texttt{id} & \texttt{INTEGER} & Clave primaria autoincremental. \\
  \texttt{message\_id} & \texttt{INTEGER} & Clave foránea indexada hacia \texttt{reported\_message}. \\
  \texttt{station} & \texttt{TEXT} & Estación que reportó: \texttt{paciente} o \texttt{recepcion}. \\
  \texttt{reported\_at} & \texttt{TEXT} & Marca temporal del reporte. \\
  \midrule
  \multicolumn{3}{l}{\textbf{Tabla \texttt{audit\_event}}} \\
  \texttt{id} & \texttt{INTEGER} & Clave primaria autoincremental. \\
  \texttt{type} & \texttt{TEXT} & \texttt{prediction}, \texttt{sentence}, \texttt{report}, \texttt{consent}, \texttt{translation\_displayed}, \texttt{reply\_sent} o \texttt{fin\_atencion}; indexado junto con \texttt{occurred\_at}. \\
  \texttt{occurred\_at} & \texttt{TEXT} & Marca temporal de ocurrencia. \\
  \texttt{station} & \texttt{TEXT} & Nulo, \texttt{paciente} o \texttt{recepcion}. \\
  \texttt{consent} & \texttt{TEXT} & Nulo, \texttt{accepted} o \texttt{declined} (RF-03.1). \\
  \texttt{lstm\_ms}, \texttt{translation\_ms}, \texttt{tts\_ms}, \texttt{total\_ms} & \texttt{REAL} & Tiempos de cada etapa en el backend, en milisegundos. \\
  \texttt{e2e\_ms} & \texttt{REAL} & Latencia extremo a extremo percibida en el frontend, entre 0 y 60.000~ms. \\
  \texttt{confidence} & \texttt{REAL} & Confianza devuelta por el modelo, en porcentaje. \\
  \texttt{source} & \texttt{TEXT} & Origen de la traducción: plantilla, caché o modelo de lenguaje. \\
  \texttt{word\_count} & \texttt{INTEGER} & Cantidad de palabras de la frase procesada. \\
  \bottomrule
\end{longtable}

\section{Diagramas de colaboración complementarios}\label{anx:colaboracion}

Se incluyen los siete diagramas de colaboración que no figuran en el cuerpo del informe. La figura~\ref{fig:dc02} corresponde al cierre de la oración y su traducción (SO-02, CU-02); la figura~\ref{fig:dc03}, a la síntesis de voz (SO-03, CU-03); la figura~\ref{fig:dc04}, al envío de una respuesta del personal (SO-04, CU-04); la figura~\ref{fig:dc05}, a la captura de una muestra de entrenamiento (SO-05, CU-06); la figura~\ref{fig:dc06}, al consentimiento de cámara (SO-06, CU-09); la figura~\ref{fig:dc07}, al registro de eventos de auditoría (SO-07, CU-11) y la figura~\ref{fig:dc09}, a la limpieza de la conversación al finalizar la atención (SO-09, CU-12).

\figuraTP{08_colaboracion_dc02}{0.8\textwidth}{DC-02: colaboración para \texttt{cerrarOracionYTraducir(secuencia)}}{fig:dc02}

\figuraTP{08_colaboracion_dc03}{\textwidth}{DC-03: colaboración para \texttt{sintetizarVoz(oracion)}}{fig:dc03}

\figuraTP{08_colaboracion_dc04}{0.85\textwidth}{DC-04: colaboración para \texttt{enviarRespuesta(mensaje)}}{fig:dc04}

\figuraTP{08_colaboracion_dc05}{\textwidth}{DC-05: colaboración para \texttt{capturarMuestra(frames, etiqueta)}}{fig:dc05}

\figuraTP{08_colaboracion_dc06}{\textwidth}{DC-06: colaboración para \texttt{otorgarConsentimiento(paciente)}}{fig:dc06}

\figuraTP{08_colaboracion_dc07}{0.9\textwidth}{DC-07: colaboración para \texttt{registrarEventoAuditoria(operacion)}}{fig:dc07}

\figuraTP{08_colaboracion_dc09}{\textwidth}{DC-09: colaboración para \texttt{limpiarConversacion()}}{fig:dc09}

\section{Registro de defectos}\label{anx:defectos}

La tabla~\ref{tab:defectos} lista los 28 defectos de la prueba de sistema: DM-LSA-001 a DM-LSA-019 se detectaron en el ciclo~1 y DM-LSA-020 a DM-LSA-028 en el ciclo~2. Al finalizar el ciclo~2 los 28 estaban cerrados. La severidad se califica de 1 (crítico) a 4 (bajo).

\begin{longtable}{>{\raggedright\arraybackslash}p{2.2cm} >{\centering\arraybackslash}p{1.0cm} >{\raggedright\arraybackslash}p{2.4cm} >{\raggedright\arraybackslash}p{8.7cm}}
  \caption{Defectos registrados en los ciclos 1 y 2}\label{tab:defectos}\\
  \toprule
  \textbf{ID} & \textbf{Sev.} & \textbf{Caso de uso o requisito} & \textbf{Título} \\
  \midrule
  \endfirsthead
  \toprule
  \textbf{ID} & \textbf{Sev.} & \textbf{Caso de uso o requisito} & \textbf{Título} \\
  \midrule
  \endhead
  DM-LSA-001 & 1 & CU-02 & La traducción agrega contenido clínico no señado e incluye fragmentos de instrucciones internas. \\
  DM-LSA-002 & 2 & CU-02 & Una traducción no clara no se descarta y se muestra al personal como una oración. \\
  DM-LSA-003 & 2 & RNF-03.1, RNF-T-06 & El backend permite cualquier origen (CORS abierto) y acepta WebSocket de orígenes ajenos. \\
  DM-LSA-004 & 2 & CU-06 & Los endpoints de captura y predicción no validan la forma de los datos: archivos corruptos, errores 500 y valores NaN. \\
  DM-LSA-005 & 2 & CU-04, CU-11 & Al enviar una respuesta del personal no se registra el evento de auditoría. \\
  DM-LSA-006 & 2 & CU-12 & No existe el botón ``Finalizar atención'' en la pantalla de recepción. \\
  DM-LSA-007 & 2 & CU-09 & No se informa al paciente cuando se deniega el acceso a la cámara. \\
  DM-LSA-008 & 2 & CU-06, CU-07 & El Modo Captura informa ``¡Guardado!'' con etiquetas o repeticiones inválidas sin haber guardado nada. \\
  DM-LSA-009 & 2 & RNF-T-02 & La ruta de landmarks por defecto depende de Internet y la cámara no se activa sin conexión. \\
  DM-LSA-010 & 2 & RNF-01.1 & La traducción con modelo de lenguaje no cumple el p95 de 1~s. \\
  DM-LSA-011 & 2 & RNF-T-05 & La interfaz presenta violaciones de accesibilidad graves o críticas. \\
  DM-LSA-012 & 3 & CU-06 & El botón ``Detener'' no aborta la grabación ni la secuencia de repeticiones. \\
  DM-LSA-013 & 3 & CU-05 & El dictado por voz sobrescribe el texto escrito y el botón de detener actúa sobre otra instancia. \\
  DM-LSA-014 & 3 & CU-05 & El dictado continuo pierde los segmentos anteriores de la transcripción. \\
  DM-LSA-015 & 3 & CU-11 & La consulta de métricas rechaza marcas de tiempo ISO~8601 con sufijo ``Z''. \\
  DM-LSA-016 & 3 & RNF-01.1 & La ruta con plantilla supera la meta de latencia en el entorno de prueba (dependiente del entorno). \\
  DM-LSA-017 & 3 & CU-10, CU-11 & La validación de los endpoints de reporte y de auditoría acepta tipos incorrectos. \\
  DM-LSA-018 & 3 & CU-01 & El protocolo WebSocket binario v2 es inalcanzable. \\
  DM-LSA-019 & 4 & CU-05 & Sin soporte de reconocimiento de voz se muestran dos avisos duplicados. \\
  DM-LSA-020 & 2 & CU-06 & En Modo Captura las muestras se guardaban vacías (sin pose, manos ni rostro). \\
  DM-LSA-021 & 1 & CU-01 & Sin ninguna persona frente a la cámara se reconocía ``dolor'' con 80~\% de confianza. \\
  DM-LSA-022 & 1 & CU-01 & Al iniciar la cámara el reposo se reconocía como ``dolor'' (96~\%). \\
  DM-LSA-023 & 2 & CU-01 & La oración se cerraba durante la pausa entre dos señas. \\
  DM-LSA-024 & 2 & CU-01 & El reconocimiento confundía ``dolor'' con ``yo''. \\
  DM-LSA-025 & 2 & RNF-04 & Con WebGL por software la detección caía a unos 1,4~fps y congelaba la interfaz. \\
  DM-LSA-026 & 2 & CU-09 & Al denegar la cámara, la carga de los modelos ocultaba el aviso de error. \\
  DM-LSA-027 & 3 & CU-06 & Durante la grabación la detección caía de 20 a unos 8~fps por efectos visuales animados. \\
  DM-LSA-028 & 4 & CU-01 & En la pantalla del paciente el panel de la cámara se agrandaba con la conversación. \\
  \bottomrule
\end{longtable}


\end{document}
